\documentclass[12pt]{article}

\usepackage[a4paper,margin=1in]{geometry}
\usepackage{setspace}
\usepackage{graphicx}
\usepackage{hyperref}
\usepackage{amsmath}
\usepackage{amssymb}
\usepackage{cite}
\usepackage{xcolor}
\usepackage{booktabs}
\usepackage{array}
\usepackage{multirow}
\usepackage{float}
\usepackage{listings}
\usepackage{algorithm}
\usepackage{algpseudocode}
\usepackage{tikz}
\usepackage{verbatim}

\setstretch{1.2}

\lstset{
    basicstyle=\ttfamily\small,
    breaklines=true,
    frame=single,
    showstringspaces=false,
}

\begin{document}

	%-------------------- COVER PAGE --------------------%

	\begin{titlepage}
		\centering

		\vspace*{1cm}
		\rule{\textwidth}{1pt}

		\vspace{2cm}

		{\Huge \textbf{Software Requirements Specification}}

		\vspace{1cm}

		{\Large for}

		\vspace{0.5cm}

		{\LARGE \textbf{SPEECH-BASED DEPRESSION SEVERITY PREDICTION SYSTEM}}

		\vspace{0.5cm}

		{\large \textit{PHQ-8 Regression using DAIC-WOZ Dataset}}

		\vspace{2cm}

		Version 2.1 approved

		\vspace{2cm}

		\textbf{Prepared by}

		\vspace{0.5cm}

		Gummadavelli Sai Balaji Koushik - 24BD1A052M \\
		Polisetty Abhishek - 24BD1A053H \\
		Rishikesh Mane - 24BD1A053P \\
		Sheelam Sai Krishna - 24BD1A053U \\
		Ageer Varun - 24BD1A05C2 \\

		\vspace{2cm}

		\textbf{KESHAV MEMORIAL INSTITUTE OF TECHNOLOGY}

		\vfill

		Date Created: \today

		\vspace{1cm}

		\rule{\textwidth}{1pt}

	\end{titlepage}

	%-------------------- REVISION HISTORY --------------------%

	\section*{Revision History}

	\begin{table}[h!]
		\centering
		\begin{tabular}{|l|l|p{6cm}|l|}
			\hline
			\textbf{Name} & \textbf{Date} & \textbf{Reason for Changes} & \textbf{Version} \\
			\hline
			G.Sai Balaji Koushik & 10-03-2026 & Initial draft of SRS with project overview and basic structure. & 1.0 \\ \hline
			Polisetty Abhishek & 11-03-2026 & Added functional requirements and speech input processing specifications. & 1.1 \\ \hline
			Rishikesh Mane & 12-03-2026 & Included feature extraction, model architecture, and ML pipeline details. & 1.2 \\ \hline
			Sheelam Sai Krishna & 13-03-2026 & Added non-functional requirements, performance metrics, and system constraints. & 1.3 \\ \hline
			Ageer Varun & 14-03-2026 & Enhanced with DAIC-WOZ dataset details, LOSO validation, and appendices. & 1.4 \\ \hline
			All Team & 27-03-2026 & Major revision: integrated multimodal fusion, transformer architecture, and technical workflows. & 2.0 \\ \hline
			G.Sai Balaji Koushik & 23-04-2026 & Performance engineering: frontend latency optimizations (response caching, request deduplication, CORS elimination, parallel API calls), backend batched metrics middleware, Swagger UI professional redesign, Vite build optimization, updated appendices. Added Section~8.8 Performance Engineering. & 2.1 \\ \hline
		\end{tabular}
	\end{table}
	\newpage

	%-------------------- ABSTRACT --------------------%

	\section*{Abstract}

	\noindent
	This Software Requirements Specification (SRS) defines a comprehensive system for automated depression severity prediction from clinical interview recordings. The system addresses a critical healthcare need: depression is a prevalent mental health disorder affecting millions globally, yet current diagnostic methods rely heavily on time-consuming subjective clinical interviews and patient self-report questionnaires. The proposed system provides objective, automated assessments by analyzing speech characteristics through multimodal machine learning, offering potential to reduce diagnostic burden and improve accessibility to mental health screening.

	\vspace{0.8em}
	\noindent
	The system builds directly upon foundational research by Alhanai et al. (2018), published in Interspeech 2018 under the title ``Detecting Depression with Audio/Text Sequence Modeling of Interviews.'' That landmark work demonstrated the effectiveness of sequence modeling for depression detection using multimodal audio and text data from the DAIC-WOZ dataset. The present specification extends this approach with contemporary deep learning advances, integrating three complementary modalities: (1) self-supervised learning (SSL) embeddings from pre-trained models (Wav2Vec2, WavLM) that capture sophisticated acoustic-phonetic patterns from raw audio, (2) linguistic representations derived from RoBERTa transformers applied to interview transcripts, capturing depression-related language characteristics, and (3) COVAREP acoustic features (74-dimensional) characterizing voice quality parameters including fundamental frequency, spectral envelope, and prosodic measures.

	\vspace{0.8em}
	\noindent
	The system implements a complete end-to-end machine learning pipeline: raw interview data (audio WAV files, speech transcripts, pre-computed acoustic features) undergoes preprocessing and validation; multimodal features are extracted in parallel from pre-trained models without early pooling to preserve temporal structure; features are normalized per-subject per-modality to account for individual speaker variation; data is organized using rigorous Leave-One-Subject-Out (LOSO) cross-validation preventing any subject from appearing in both training and test sets; the hybrid transformer-BiLSTM model with attention-based multimodal fusion is trained for each fold; and finally comprehensive evaluation metrics (R², MAE, RMSE) are computed on held-out test subjects. The neural architecture combines transformer encoders (capturing long-range acoustic-linguistic dependencies via self-attention), bidirectional LSTMs (capturing local temporal patterns and sequential dynamics), and learned attention-based multimodal fusion (enabling dynamic inter-modal weighting) to produce continuous PHQ-8 predictions in the clinically relevant range 0--24.

	\vspace{0.8em}
	\noindent
	Evaluation is conducted using LOSO cross-validation across approximately 189 subjects from the DAIC-WOZ dataset, the same benchmark employed in Alhanai et al., ensuring subject-level generalization and preventing optimistic performance overestimation. The system targets regression performance of $R^2 > 0.50$ with mean absolute error below 3.0 PHQ-8 points, establishing a competitive computational baseline for automated depression severity assessment. This SRS comprehensively specifies functional requirements (data ingestion, feature extraction, model training, evaluation), non-functional requirements (performance targets, scalability, reliability, security), detailed system architecture with modular design, complete workflow descriptions with algorithmic specifications, and comparative justification tables for architectural and methodological decisions grounded in contemporary research and practical constraints.

	\vspace{1em}
	\noindent
	\textbf{Keywords:} depression detection, speech analysis, PHQ-8 regression, multimodal machine learning, LOSO cross-validation, transformer architectures, acoustic features, sequence modeling, clinical interview analysis.

	\newpage

	%-------------------- TABLE OF CONTENTS --------------------%

	\pagenumbering{roman}

	\tableofcontents

	\newpage

	\pagenumbering{arabic}
	\setcounter{page}{1}

	% =====================================================================
	% 1. INTRODUCTION
	% =====================================================================

	\section{Introduction}

	This document provides a comprehensive Software Requirements Specification (SRS) for the \textit{Speech-Based Depression Severity Prediction System}. The system is designed to analyze speech recordings from clinical interviews and predict PHQ-8 (Patient Health Questionnaire-8) depression severity scores using multimodal machine learning techniques. This SRS defines functional requirements, non-functional requirements, system architecture, and development constraints.

	\subsection{Purpose}

	The primary purpose of this SRS document is to:

	\begin{itemize}
		\item Define functional and non-functional requirements for the depression severity prediction system
		\item Specify the technical architecture combining audio, text, and acoustic modalities
		\item Document the machine learning pipeline including feature extraction and model training
		\item Establish evaluation criteria and validation methodology (LOSO cross-validation)
		\item Serve as a reference for developers, researchers, and project supervisors
	\end{itemize}

	The depression severity prediction system addresses the need for objective, automated assessment of mental health status through speech analysis. Unlike traditional clinical approaches that rely on subjective interviews, this system provides quantitative predictions grounded in acoustic and linguistic speech characteristics.

	\subsection{Document Conventions}

	Throughout this document, the following conventions are employed:

	\begin{itemize}
		\item \textbf{Functional Requirements (FR):} Specific capabilities the system must provide
		\item \textbf{Non-Functional Requirements (NFR):} Performance, reliability, and quality attributes
		\item \textbf{Priority Levels:} High (critical), Medium (important), Low (optional)
		\item \textbf{Modalities:} Audio (speech signal), Text (transcripts), Acoustic (COVAREP features)
		\item \textbf{LOSO:} Leave-One-Subject-Out cross-validation for subject-level generalization
	\end{itemize}

	\subsection{Intended Audience}

	This document is intended for multiple stakeholder groups:

	\begin{itemize}
		\item \textbf{Machine Learning Engineers:} To understand system architecture, model design, and implementation details
		\item \textbf{Project Supervisors:} To review requirements, architecture decisions, and validation methodology
		\item \textbf{Researchers:} To understand the speech-based mental health analysis approach and experimental design
		\item \textbf{Future Maintainers:} To reference system design during updates and enhancements
		\item \textbf{Clinical Evaluators:} To assess system performance and clinical validity
	\end{itemize}

	\subsection{Product Scope}

	The speech-based depression severity prediction system focuses on automated analysis of clinical interview recordings. The scope includes:

	\begin{itemize}
		\item \textbf{Input Processing:} Audio recordings (WAV 16kHz), speech transcripts, acoustic feature matrices (COVAREP)
		\item \textbf{Feature Engineering:} Self-supervised learning models (Wav2Vec2, WavLM), transformer language models (RoBERTa), traditional acoustic features (MFCC)
		\item \textbf{Model Architecture:} Hybrid transformer-BiLSTM with attention-based multimodal fusion
		\item \textbf{Output:} Continuous PHQ-8 score predictions (range 0--24)
		\item \textbf{Validation:} LOSO cross-validation across ~189 subjects
	\end{itemize}

	\textbf{Out of Scope:} Real-time inference optimization, mobile deployment, clinical decision support integration (Phase 2), multimodal video analysis (facial/body language).

	\subsection{Key References}

	\begin{itemize}
		\item Alhanai, T., Ghassemi, M. M., \& Glass, J. R. (2018). ``Detecting depression with audio/text sequence modeling of interviews.'' In \textit{Interspeech 2018}.
		\item DAIC-WOZ Dataset – University of Southern California, Informatics and Data-Driven Culture Lab
		\item Vaswani, A., et al. (2017). ``Attention is All You Need.'' In \textit{NIPS 2017}.
		\item PyTorch and Hugging Face Transformers Documentation
	\end{itemize}

	% =====================================================================
	% 2. OVERALL DESCRIPTION
	% =====================================================================

	\section{Overall Description}

	\subsection{Product Perspective}

	The speech-based depression severity prediction system represents a bridge between clinical mental health assessment and computational speech analysis. From a high-level perspective, the system operates as follows:

	\begin{enumerate}
		\item Patient participates in a structured clinical interview
		\item Interview audio and transcript are collected
		\item System processes multimodal data through feature extraction pipeline
		\item Machine learning model predicts PHQ-8 severity score
		\item Results are presented with confidence estimates and visualization
	\end{enumerate}

	The system is positioned as a \textbf{research tool} for validating automated depression assessment rather than a direct clinical diagnostic instrument. Future iterations may enable integration into telemedicine platforms or clinical decision support systems.

	Key design characteristics:
	\begin{itemize}
		\item \textbf{Modular Architecture:} Audio processing, text encoding, and acoustic feature extraction can be updated independently
		\item \textbf{Transfer Learning Based:} Leverages pre-trained models (Wav2Vec2, RoBERTa) to maximize performance on limited data
		\item \textbf{Multimodal Integration:} Attention mechanisms dynamically weight contributions from different modalities
		\item \textbf{Subject-Level Validation:} LOSO cross-validation prevents data leakage and ensures honest generalization estimates
	\end{itemize}

	\subsection{Product Functions}

	The system implements the following core functions:

	\begin{enumerate}
		\item \textbf{Audio Feature Extraction:}
		\begin{itemize}
			\item Load WAV files and resample to 16 kHz mono
			\item Extract self-supervised learning embeddings (Wav2Vec2: 768-dim, WavLM: 768-dim)
			\item Extract MFCC features with delta and delta-delta derivatives
			\item Chunk audio into overlapping 3-second segments with 2-second stride
			\item Maintain temporal sequence structure (no early averaging)
		\end{itemize}

		\item \textbf{Text Feature Extraction:}
		\begin{itemize}
			\item Clean and normalize speech transcripts
			\item Encode transcripts using RoBERTa transformer (768-dim CLS token)
			\item Handle both manual and ASR-generated transcripts
		\end{itemize}

		\item \textbf{Acoustic Feature Processing:}
		\begin{itemize}
			\item Load pre-extracted COVAREP features (74-dimensional)
			\item Aggregate and normalize acoustic characteristics
		\end{itemize}

		\item \textbf{Multimodal Fusion:}
		\begin{itemize}
			\item Combine audio, text, and acoustic modalities using attention mechanisms
			\item Project each modality to common 256-dim representation
			\item Apply multi-head attention over modality tokens
			\item Generate fused embedding vector
		\end{itemize}

		\item \textbf{PHQ-8 Prediction:}
		\begin{itemize}
			\item Process fused embedding through dense regression head
			\item Output continuous score predictions (0--24 range)
			\item Perform LOSO cross-validation evaluation
		\end{itemize}

		\item \textbf{Results Management:}
		\begin{itemize}
			\item Log training metrics and model checkpoints
			\item Compute R², MAE, RMSE evaluation metrics
			\item Generate visualizations and reports
		\end{itemize}
	\end{enumerate}

	\subsection{User Classes and Characteristics}

	\begin{description}
		\item[ML Researchers] Users with advanced Python and PyTorch expertise who conduct model development, hyperparameter tuning, and algorithm innovation. Requirements: full access to codebase, detailed logging, configuration flexibility.

		\item[Clinical Researchers] Users who validate system performance against clinical benchmarks. Requirements: evaluation metrics, subject-level breakdowns, statistical reporting.

		\item[System Evaluators] Quality assurance personnel testing reproducibility and robustness. Requirements: deterministic execution, comprehensive test cases, performance profiling.

		\item[API Consumers] External systems or applications using inference API for depression prediction. Requirements: standardized REST interface, latency guarantees, error handling.
	\end{description}

	\subsection{Operating Environment}

	The system operates within the following technical environment:

	\begin{itemize}
		\item \textbf{Hardware:} GPU-accelerated compute (NVIDIA CUDA 11.x+, 8GB+ VRAM for training)
		\item \textbf{Software:} Python 3.8+, PyTorch 1.10+, Transformers 4.15+, librosa 0.9+
		\item \textbf{Data Format:} WAV audio, CSV transcripts/features, JSON configuration and results
		\item \textbf{Deployment:} Linux/Unix workstations, Docker containerization, cloud GPU instances (optional)
	\end{itemize}

	\subsection{Design and Implementation Constraints}

	Critical constraints guiding system design:

	\begin{enumerate}
		\item \textbf{Temporal Structure Preservation:} Audio sequences must maintain temporal information throughout the pipeline. No early averaging or pooling before temporal modeling layers.

		\item \textbf{Subject-Level Isolation:} LOSO cross-validation requires complete subject separation. No training data from test subjects can appear in development sets.

		\item \textbf{Per-Speaker Normalization:} Feature normalization must be applied per-subject to account for speaker-specific variation without biasing the model.

		\item \textbf{Attention-Based Fusion Required:} Multimodal integration must use attention mechanisms, not simple concatenation, to enable learned inter-modal weighting.

		\item \textbf{Lightweight Transformer:} Transformer encoder limited to 2 layers and 4--8 attention heads to prevent overfitting on ~189 subjects.

		\item \textbf{No PHQ Sub-Score Leakage:} Only total PHQ-8 score used for training; individual symptom dimensions excluded to prevent information leakage.
	\end{enumerate}

	\subsection{Assumptions and Dependencies}

	\begin{itemize}
		\item \textbf{Dataset Availability:} Access to DAIC-WOZ dataset (~189 complete subjects with audio, transcripts, COVAREP, and PHQ-8 labels)
		\item \textbf{Pre-trained Models:} Availability of Wav2Vec2, WavLM, and RoBERTa via Hugging Face model hub
		\item \textbf{GPU Infrastructure:} Sufficient GPU resources for feature extraction and model training
		\item \textbf{Stable Dependencies:} PyTorch and related libraries maintain API compatibility during development
		\item \textbf{Data Privacy:} DAIC-WOZ data appropriately anonymized with informed consent obtained
	\end{itemize}

	% =====================================================================
	% 3. SYSTEM ARCHITECTURE & DATA FLOW
	% =====================================================================

	\section{System Architecture and Data Flow}

	\subsection{Non-Technical System Architecture}

	\begin{figure}[h]
		\centering
		\includegraphics[width=0.9\linewidth, height=0.4\textheight, keepaspectratio]{system architecture non technical.png}
		\caption{Non-Technical System Architecture: Complete data flow pipeline from interview data collection through preprocessing, parallel feature extraction from three modalities, LOSO cross-validation training, to final depression severity predictions. Shows subject isolation strategy and modular processing stages.}
		\label{fig:system_arch_nontechnical}
	\end{figure}

	\begin{enumerate}

    \item \textbf{Input Acquisition}  
    The system accepts three complementary input modalities derived from clinical interview recordings:
    \begin{itemize}
        \item \textbf{Audio Data:} Raw speech signals recorded at 16 kHz.
        \item \textbf{Transcripts:} Textual representation of the interview obtained through manual transcription or automatic speech recognition.
        \item \textbf{Acoustic Features:} Pre-computed COVAREP feature matrices representing voice production characteristics.
    \end{itemize}

    \item \textbf{Data Validation}  
    Upon receiving the inputs, the system performs integrity checks to ensure:
    \begin{itemize}
        \item Files are readable and accessible.
        \item No corruption is present in the data.
        \item All required fields and modalities are available.
    \end{itemize}

    \item \textbf{Parallel Feature Processing}  
    The validated multimodal data is processed through a parallel architecture where each modality is handled independently:
    \begin{itemize}
        \item \textbf{Audio Processing:} Audio signals are segmented into overlapping chunks and processed using self-supervised learning models along with traditional acoustic feature extractors.
        \item \textbf{Text Processing:} Transcripts are encoded using transformer-based language models to capture semantic and linguistic patterns.
        \item \textbf{Acoustic Feature Handling:} COVAREP features are loaded and normalized.
    \end{itemize}
    This parallel processing design improves computational efficiency and reduces overall processing time.

    \item \textbf{Feature Caching}  
    Extracted features from all modalities are stored in binary format on disk. This caching mechanism avoids redundant computations during subsequent training iterations, thereby significantly improving system efficiency.

    \item \textbf{Dataset Organization and Splitting}  
    The system organizes data using a \textbf{Leave-One-Subject-Out (LOSO) cross-validation} strategy:
    \begin{itemize}
        \item For each iteration, one subject is designated as the test set.
        \item The remaining subjects are divided into training and validation sets.
        \item This ensures no subject appears in multiple splits, maintaining strict data separation.
    \end{itemize}

    \item \textbf{Model Training and Prediction}  
    For each LOSO fold:
    \begin{itemize}
        \item The neural network is trained using the training set.
        \item Model performance is validated using the validation set.
        \item Predictions are generated for the held-out test subject.
    \end{itemize}

    \item \textbf{Performance Aggregation and Reporting}  
    After completing all folds:
    \begin{itemize}
        \item Performance metrics such as $R^2$, MAE, and RMSE are computed for each fold.
        \item Cross-fold mean and standard deviation statistics are calculated.
        \item Results are compiled into comprehensive reports, including visualizations, error analysis, and outputs in both JSON and human-readable formats.
    \end{itemize}

\end{enumerate}

	\subsection{Technical System Architecture}

	\begin{figure}[h]
		\centering
		\includegraphics[width=0.9\linewidth, height=0.5\textheight, keepaspectratio]{system architecture_technical.png}
		\caption{Technical System Architecture: Detailed neural network layer specification showing tensor dimensions at each transformation. Depicts transformer encoder plus bidirectional LSTM sequence processing, attention-based temporal pooling, multimodal fusion combining audio/text/acoustic modalities, and dense regression head outputting continuous PHQ-8 predictions.}
		\label{fig:system_arch_technical}
	\end{figure}

	The technical system architecture provides detailed layer-by-layer specifications of the neural network processing pipeline and internal data transformations. The system comprises five major functional components: 
	\begin{enumerate}
    \item \textbf{Data Ingestion Module}  
    Accepts multimodal inputs and applies padding and masking to handle variable-length sequences. Audio sequences may range from 100 to 2000+ timesteps depending on interview duration. Text embeddings are represented as fixed 768-dimensional vectors, while acoustic features are consistently 74-dimensional.

    \item \textbf{Preprocessing Module}  
    Performs audio resampling to a standardized 16 kHz sampling rate. Transcript preprocessing includes XML tag removal and punctuation normalization. Acoustic feature validation is carried out to ensure the absence of missing or corrupted values.

    \item \textbf{Feature Extraction Module}  
    Executes three parallel processing pipelines:
    \begin{itemize}
        \item \textbf{Audio Features:} Generates temporal sequences $(T \times 900)$ by concatenating Wav2Vec2 embeddings (768-dimensional) with aligned MFCC features (39--120 dimensions).
        \item \textbf{Text Features:} Produces RoBERTa CLS token embeddings (768-dimensional), capturing semantic sentence-level representations.
        \item \textbf{Acoustic Features:} Aggregates COVAREP features into a 74-dimensional vector.
    \end{itemize}
    Per-subject normalization is applied independently to each modality using a StandardScaler fitted exclusively on training data.

    \item \textbf{Neural Network Module}  
    Implements a hybrid deep learning architecture:
    \begin{itemize}
        \item \textbf{Transformer Encoder:} Consists of 2 layers with 4--8 attention heads and hidden dimensions of 256--384. It captures long-range dependencies in audio sequences, enabling identification of depression-related prosodic patterns such as monotonic speech and reduced pitch variation.
        
        \item \textbf{Bidirectional LSTM:} Includes 1--2 layers with 256 hidden units per direction (512-dimensional combined output). It models local temporal dependencies not fully captured by the transformer.
        
        \item \textbf{Attention-based Temporal Pooling:} Assigns learned weights to each timestep, emphasizing relevant segments while reducing the influence of silence or padded regions.
        
        \item \textbf{Multimodal Fusion:} Each modality is projected into a 256-dimensional space using linear layers with ReLU activation and layer normalization. These are treated as tokens and processed using multi-head attention (4 heads). Final attention pooling produces a unified 256-dimensional representation.
    \end{itemize}

    \item \textbf{Regression Head Module}  
    Processes the fused representation through fully connected layers:
    \begin{itemize}
        \item $256 \rightarrow 256$ with ReLU activation and 0.3 dropout
        \item $256 \rightarrow 64$ with ReLU activation
        \item $64 \rightarrow 1$ producing continuous output
    \end{itemize}
    The output is clipped to the valid PHQ-8 score range $[0, 24]$.

    \item \textbf{Regularization and Optimization Techniques}  
    \begin{itemize}
        \item Padding masks prevent attention over invalid padded positions.
        \item Layer normalization stabilizes gradient flow across layers.
        \item Dropout is applied to mitigate overfitting given the limited dataset size (189 subjects).
    \end{itemize}

\end{enumerate}

	\subsection{System Modules and Data Flow}

	The system is organized into modular components enabling independent updates and testing:

	\begin{description}
		\item[Data Input Module] Manages access to DAIC-WOZ dataset files and validates data integrity. Supports batch loading of audio (WAV format), transcripts (plaintext/CSV), and acoustic features (CSV). Reports data quality issues and incomplete samples.

		\item[Preprocessing Module] Performs text cleaning through whitespace normalization and punctuation standardization, audio validation checking for header integrity and non-zero content, and missing data detection identifying subjects with incomplete modalities. Generates preprocessing reports for transparency.

		\item[Feature Extraction Module] Parallelizes extraction across three modalities without sequential dependencies. Audio extraction yields temporal sequences (T×900) from SSL+MFCC concatenation. Text extraction yields pooled 768-dim CLS tokens from RoBERTa. Acoustic extraction loads and aggregates COVAREP (74-dim). All features are cached to disk in PyTorch tensor format.

		\item[Dataset Module] Implements PyTorch Dataset and DataLoader with custom collate functions handling variable-length sequences, padding to maximum length, attention mask generation, and LOSO subject-wise splitting preventing any subject appearing in multiple splits.

		\item[Model Architecture Module] Defines the complete neural network implementing transformer encoder, bidirectional LSTM, attention pooling, multimodal fusion (audio+text+COVAREP), and regression head. Supports configuration-based hyperparameter adjustment.

		\item[Training Module] Implements the complete LOSO cross-validation loop: for each subject, initializes a fresh model, trains on remaining subjects with validation monitoring, applies early stopping and learning rate scheduling, saves best checkpoints, and evaluates on held-out test subject.

		\item[Evaluation Module] Computes comprehensive metrics (R², MAE, RMSE) on test and validation sets, generates per-subject error analysis, creates visualizations (scatter plots, residual histograms), and aggregates results across folds with statistical summaries.

		\item[Inference API Module] Provides REST API endpoints accepting audio and transcript inputs, orchestrating complete feature extraction and prediction pipeline, returning JSON responses with PHQ-8 predictions and optional confidence intervals.
	\end{description}

	\subsection{Feature Extraction Pipeline}

	The feature extraction pipeline processes multimodal data in parallel:

	\subsubsection{Audio Feature Extraction}

	Audio processing preserves temporal structure:

	\begin{itemize}
		\item Load WAV file, resample to 16 kHz mono
		\item Chunk audio into overlapping segments (3-second chunks, 2-second stride)
		\item For each chunk:
		\begin{itemize}
			\item Extract Wav2Vec2 embeddings (768-dim, temporal sequence preserved)
			\item Extract WavLM embeddings (768-dim, temporal sequence preserved)
			\item Extract MFCC (13--40 coefficients + delta + delta-delta)
			\item Align temporal dimensions via interpolation
			\item Concatenate SSL + MFCC features $\approx 900$ dimensions
		\end{itemize}
		\item Aggregate chunks into subject-level sequence (T $\times$ 900 tensor)
		\item Apply per-subject StandardScaler normalization
	\end{itemize}

	\subsubsection{Text Feature Extraction}

	Transcript encoding is performed once per subject:

	\begin{itemize}
		\item Clean transcript (remove XML, normalize whitespace)
		\item Encode using RoBERTa-base through Hugging Face transformers
		\item Extract CLS token embedding (768-dim)
		\item Store as fixed-dimensional vector per subject
	\end{itemize}

	\subsubsection{Acoustic Feature Extraction}

	COVAREP features are aggregated from provided matrices:

	\begin{itemize}
		\item Load COVAREP CSV (74-dimensional acoustic features)
		\item Perform mean aggregation across frames if frame-level data
		\item Store as 74-dim vector per subject
	\end{itemize}

	% Workflow diagram placeholder for workflow_technical
	% NOTE: TikZ code for technical workflow diagram
	% LOSO Cross-Validation Pipeline

	\subsection{Workflow Description and Pipeline}

	\begin{figure}[h]
		\centering
		\includegraphics[width=0.9\linewidth, height=0.45\textheight, keepaspectratio]{workflow.png}
		\caption{Technical LOSO Cross-Validation Workflow: Shows complete iterative pipeline for each subject including parallel feature extraction, per-subject normalization, model training loop with early stopping and learning rate scheduling, test evaluation, and cross-fold aggregation of R², MAE, RMSE metrics.}
		\label{fig:workflow_technical}
	\end{figure}

	The workflow pipeline orchestrates the complete end-to-end process from raw interview data to final depression severity predictions using rigorous Leave-One-Subject-Out (LOSO) cross-validation. This validation strategy is critical because it ensures that when evaluating model performance on a held-out test subject, neither that subject's data nor any subject sharing demographic or clinical characteristics appears in the training or validation sets, preventing overly optimistic performance estimates and providing honest subject-level generalization metrics. The workflow implements the following sequential operational stages:

	\begin{enumerate}
		\item \textbf{Data Loading:} Load complete dataset metadata identifying all available subjects and corresponding data file paths for audio (WAV format at 16 kHz), transcripts (plaintext or CSV with speaker annotations), COVAREP acoustic features (CSV with 74 dimensions), and PHQ-8 target scores.

		\item \textbf{LOSO Split Generation:} For each of $N$ subjects in the dataset, systematically designate that subject as the exclusive test set (containing exactly 1 subject), then randomly partition the remaining $N-1$ subjects into two disjoint sets: validation set (10--15\% of non-test subjects, typically 16--24 subjects for $N=189$) and training set (85--90\% of non-test subjects, typically 155--170 subjects), ensuring a validation set sufficient for reliable monitoring yet preserving maximal training sample size.

		\item \textbf{Parallel Multimodal Feature Extraction:} For all subjects in the current fold's training and validation sets, execute feature extraction in parallel across three independent modality pipelines: (a) \textbf{Audio extraction} processes each subject's audio file, chunks it into overlapping 3-second segments with 2-second stride, extracts Wav2Vec2 embeddings (768-dimensional) and WavLM embeddings (768-dimensional) preserving full temporal sequence through the neural network without early pooling, extracts MFCC features (13--40 coefficients plus delta and delta-delta derivatives for 39--120 total dimensions), aligns temporal dimensions through interpolation, and concatenates SSL and MFCC features yielding T×900 dimensional temporal sequences; (b) \textbf{Text extraction} loads and cleans the subject's interview transcript (removing speaker labels, normalizing whitespace, standardizing punctuation), tokenizes and encodes through RoBERTa-base transformer via Hugging Face API, pools the final layer's CLS token to obtain a fixed 768-dimensional semantic representation; (c) \textbf{Acoustic extraction} loads the subject's COVAREP feature matrix from CSV, performs mean aggregation across frame dimension (if frame-level features are provided) to produce a single 74-dimensional acoustic vector, optionally extends with standard deviation statistics for temporal variability capture.

		\item \textbf{Feature Normalization:} Compute StandardScaler normalization statistics (mean and standard deviation) separately for each modality (audio, text, acoustic) using \textit{training set only}, strictly preventing information leakage from validation or test subjects. Apply these training-fitted normalization parameters to scale training, validation, and test samples identically: $x_{normalized} = (x - \mu_{train}) / \sigma_{train}$. This per-subject per-modality normalization accounts for individual speaker vocal variation without introducing systematic bias.

		\item \textbf{Feature Caching:} Serialize all normalized features to disk as PyTorch binary tensors in efficient HDF5 or PT formats, storing audio sequences in dedicated directory (data/processed/audio\_features/), text embeddings in directory (data/processed/text\_features/), and acoustic features in directory (data/processed/covarep\_features/). Caching dramatically accelerates subsequent training iterations, reducing epoch duration from minutes to seconds by eliminating expensive recomputation of SSL embeddings and RoBERTa encodings.

		\item \textbf{DataLoader Construction:} Instantiate PyTorch Dataset objects wrapping cached feature files, implement custom collate function that: (a) retrieves variable-length audio sequences from storage, (b) \textbf{pads sequences to maximum length} ($T_{max}$) observed in training set, (c) generates attention masks as binary tensors (1 for valid content, 0 for padded positions), (d) retrieves fixed-dimensional text and acoustic embeddings, (e) retrieves target PHQ-8 labels, (f) batches samples to specified batch size (8--16). Create separate DataLoaders for training (shuffle=True), validation (shuffle=False), and test (shuffle=False) to ensure proper gradient computation and reproducible evaluation.

		\item \textbf{Model Initialization:} Create a new hybrid transformer-BiLSTM-fusion model instance with random weight initialization. This fresh initialization for each LOSO fold is critical: reusing weights from prior folds would introduce spurious correlations and violate the independence assumption of cross-validation.

		\item \textbf{Training Loop Execution:} Implement 30--50 epoch training loop with the following components: (a) forward pass through complete neural architecture computing predictions; (b) MSE loss computation comparing predictions to target PHQ-8 scores; (c) backpropagation computing parameter gradients via automatic differentiation; (d) gradient clipping applying L2-norm constraint at threshold 1.0 to prevent exploding gradients; (e) parameter update via AdamW optimizer with learning rate $1e^{-4}$ and weight decay $1e^{-5}$; (f) epoch-level validation on validation set computing R² (coefficient of determination), MAE (mean absolute error), and RMSE (root mean squared error); (g) learning rate scheduling: if validation loss fails to improve for 2--3 consecutive epochs, reduce learning rate by factor 0.5; (h) checkpointing: save complete model state whenever validation $R^2$ exceeds previous best; (i) early stopping: terminate training if validation loss shows no improvement for 5 consecutive epochs, preserving the best-performing model.

		\item \textbf{Test Evaluation:} Load the best model checkpoint identified during training, perform forward pass (without gradient computation, with batch normalization in inference mode) on the single held-out test subject, compute per-fold R², MAE, and RMSE metrics quantifying prediction accuracy on this unseen individual, store results with comprehensive metadata (fold number, subject ID, prediction confidence measures).

		\item \textbf{Cross-Fold Aggregation:} After completing training and evaluation for all $N$ subjects: (a) compile per-fold R², MAE, RMSE arrays; (b) compute cross-fold mean and standard deviation for each metric, yielding final performance estimates (e.g., mean $R^2 \pm$ std); (c) compute per-subject absolute prediction errors $|y_i - \hat{y}_i|$ across all folds; (d) rank subjects by prediction error identifying easiest and hardest cases; (e) create visualization artifacts including scatter plots of predicted vs. actual PHQ-8 scores with regression line, residual histograms characterizing error distribution, per-fold performance bar charts; (f) generate comprehensive JSON results file including all metrics, fold-level details, and hyperparameters; (g) output human-readable performance reports summarizing results and comparing to literature benchmarks.
	\end{enumerate}

	This workflow ensures that each subject serves exactly once as a test case, subject-level generalization is rigorously evaluated without data leakage, computational efficiency is maximized through feature caching, and results are reproducible and thoroughly documented for research publication and clinical validation.

	\subsection{Pipeline and Models: Comprehensive Overview}

	\subsubsection{End-to-End Pipeline Stages}

	\textbf{Stage 1: Audio Feature Extraction}
	\begin{itemize}
		\item \textbf{Wav2Vec2 Model} (\texttt{facebook/wav2vec2-base-960h}): Pre-trained on 960 hours LibriSpeech. Extracts 768-dimensional embeddings per 20ms frame, capturing phonetic patterns.
		\item \textbf{WavLM Model} (\texttt{microsoft/wavlm-base}): Pre-trained on 50,000+ hours multilingual speech. Extracts 768-dimensional embeddings per frame, robust to speaker variation and prosody.
		\item \textbf{MFCC Features}: 40 Mel-frequency cepstral coefficients + delta + delta-delta = 39--120 dimensions per frame. Captures spectral texture and speech dynamics.
		\item \textbf{Audio Output}: Three sources concatenated: 768 (Wav2Vec2) + 768 (WavLM) + 40--120 (MFCC) = T×900 temporal sequences.
	\end{itemize}

	\textbf{Stage 2: Text Feature Extraction}
	\begin{itemize}
		\item \textbf{RoBERTa Model} (\texttt{facebook/roberta-base}): 12 transformer layers, 12 heads, 768-dimensional embeddings. Pre-trained on 160GB diverse English text (CommonCrawl, Wikipedia, Reddit).
		\item \textbf{Pooling Strategy}: [CLS] token extraction (special token designed to aggregate document meaning).
		\item \textbf{Text Output}: Fixed 768-dimensional vector per subject capturing semantic interview content.
	\end{itemize}

	\textbf{Stage 3: Acoustic Features}
	\begin{itemize}
		\item \textbf{COVAREP}: Computational Vocal Resonance Period—74-dimensional feature vectors capturing voice quality, fundamental frequency F0, spectral envelope, and voice stability (HNR, shimmer, jitter).
		\item \textbf{Aggregation}: Mean over all frames yields single 74-dimensional vector per subject.
		\item \textbf{Acoustic Output}: Fixed 74-dimensional vector per subject representing average voice characteristics.
	\end{itemize}

	\textbf{Stage 4: Multimodal Fusion}
	\begin{itemize}
		\item \textbf{Audio Processing}: T×900 → Transformer Encoder (2 layers, 4--8 heads, 256--384 hidden) → BiLSTM (1--2 layers, 256 hidden) → Attention Pooling → 512-dimensional vector.
		\item \textbf{Projection Layer}: Project audio (512), text (768), acoustic (74) to common 256-dimensional space using ReLU + LayerNorm.
		\item \textbf{Attention Fusion}: Multi-head attention (4 heads) over 3 modality tokens → 256-dimensional fused vector.
		\item \textbf{Regression Head}: 256→256 (ReLU, Dropout 0.3) → 64 (ReLU) → 1 (Linear) → Clamp to [0, 24].
	\end{itemize}

	\textbf{Stage 5: LOSO Cross-Validation}
	\begin{itemize}
		\item 189 folds (one per subject)
		\item Per fold: Test = 1 held-out subject, Validation = 20 subjects (15%), Training = 168 subjects (85%)
		\item Ensures subject-level generalization without data leakage
	\end{itemize}

	\subsubsection{Model Architecture Details}

	\textbf{Transformer Encoder Specification}
	\begin{itemize}
		\item Layers: 2 (balances capacity with overfitting prevention on 189 subjects)
		\item Attention Heads: 4--8 per layer
		\item Hidden Dimension: 256--384 (per layer width)
		\item Dropout: 0.1 (regularization)
		\item Purpose: Captures long-range temporal dependencies (depression manifests as consistent patterns across interviews)
	\end{itemize}

	\textbf{BiLSTM Specification}
	\begin{itemize}
		\item Layers: 1--2 bidirectional
		\item Hidden Dimension: 256 per direction (512 total output per timestep)
		\item Function: Forward LSTM processes left-to-right, backward LSTM right-to-left
		\item Benefit: Captures local temporal patterns (speech rate, pause duration) and provides bidirectional context
	\end{itemize}

	\textbf{Training Hyperparameters}
	\begin{itemize}
		\item Optimizer: AdamW (adaptive learning rate with weight decay)
		\item Learning Rate: $10^{-4}$ (conservative for stability on limited data)
		\item Weight Decay: $10^{-5}$ (L2 regularization)
		\item Batch Size: 8--16 (small size adds regularization noise)
		\item Early Stopping: Patience = 5 epochs (monitor validation loss)
		\item Dropout in Head: 0.3 (stronger regularization protecting final predictions)
		\item Audio Chunking: 3-second chunks with 2-second stride (50% overlap)
	\end{itemize}

	\textbf{Model Parameter Count: ~5.2M}
	\begin{itemize}
		\item Transformer: ~2.5M
		\item BiLSTM: ~0.8M
		\item Projections + Fusion: ~0.9M
		\item Regression Head: ~0.2M
		\item Rationale: Conservative on 189 subjects (~30 samples per parameter) prevents overfitting
	\end{itemize}

	\subsubsection{Feature Types Explained}

	\textbf{Self-Supervised Learning (SSL) Embeddings}

	SSL models learn representations from unlabeled data via contrastive pretraining: masking random audio segments and predicting from context. \textbf{Wav2Vec2} and \textbf{WavLM} both output 768-dimensional vectors per 20ms frame. Wav2Vec2 excels at phonetic content; WavLM captures speaker characteristics and prosody. Combining both provides complementary perspectives.

	\textbf{MFCC (Mel-Frequency Cepstral Coefficients)}

	MFCCs model human auditory perception using logarithmic frequency scaling. 40 coefficients per frame capture spectral energy distribution. Delta features (velocity) measure frame-to-frame change; delta-delta (acceleration) measure change rate. Total: 39--120 dimensions per frame. Depression correlates with: flat MFCCs (monotone), low deltas (slow speech), low delta-deltas (reduced dynamics).

	\textbf{RoBERTa Language Model}

	12-layer transformer pre-trained on 160GB text. Bidirectional context enables words to be represented differently based on surrounding meaning. [CLS] token is designed to aggregate document meaning; extracting its 768-dimensional embedding pools the entire transcript into a single vector capturing interview semantics (word choice, sentiment, linguistic markers of depression).

	\textbf{COVAREP Acoustic Features (74-dim)}

	Expert-designed features measuring voice production: F0 value (pitch) and variation, spectral envelope (voice brightness), voice quality (HNR = harmonics-to-noise ratio, shimmer = amplitude variation, jitter = pitch variation). Depression manifests as: monotone F0 (reduced variation), high HNR/shimmer/jitter (vocal instability), duller sound (reduced spectral energy).

	\textbf{Why Three Modalities?}

	Single modality insufficient. Audio captures prosody/acoustics. Text captures semantic content. COVAREP captures voice quality. Multimodal fusion with attention learns to weight modalities: clear speech + negative language → weight text; unclear speech + normal language → weight audio/acoustic. This adaptive weighting is superior to naive concatenation.

	% =====================================================================
	% 4. DATASET & METHODOLOGY
	% =====================================================================

	\section{Dataset and Methodology}

	\subsection{DAIC-WOZ Dataset Description}

	\subsubsection{Dataset Overview}

	The DAIC-WOZ (Depression Anxiety to Help Through TeleHealth Interventions to Overcome Depression--Wizard-of-Oz) dataset provides clinical interview data collected by the University of Southern California's Informatics and Data-Driven Culture Lab. The dataset comprises:

	\begin{itemize}
		\item \textbf{Approximately 189 interview sessions} with targeted depression assessments
		\item \textbf{Audio recordings} in WAV format at 16 kHz sampling rate
		\item \textbf{Speech transcripts} with speaker labels (participant, interviewer)
		\item \textbf{COVAREP acoustic features} (74-dimensional feature vectors)
		\item \textbf{PHQ-9/PHQ-8 severity scores} ranging 0--27 (PHQ-9) or 0--24 (PHQ-8)
	\end{itemize}

	The dataset was collected under IRB approval with informed participant consent and represents a naturalistic but partially structured interview environment (virtual interviewer ``Ellie'' controlled by human operators).

	\subsubsection{Data Distribution}

	PHQ-8 score distribution across dataset:
	\begin{itemize}
		\item Minimal Depression (0--4): ~45\% of subjects
		\item Mild Depression (5--9): ~35\% of subjects
		\item Moderate Depression (10--14): ~15\% of subjects
		\item Severe Depression (15--24): ~5\% of subjects
	\end{itemize}

	Interview duration: 7--35 minutes typically. Subject diversity includes multiple age groups, genders, and demographic backgrounds.

	\subsubsection{Feature Modalities}

	Three complementary modalities are available:

	\begin{enumerate}
		\item \textbf{Audio Waveforms:} Raw speech signal enabling extraction of prosodic, spectral, and self-supervised learning features

		\item \textbf{COVAREP Acoustic Features:} 74-dimensional pre-computed features capturing:
		\begin{itemize}
			\item Fundamental frequency and voicing characteristics
			\item Spectral envelope coefficients
			\item Voice quality measures
			\item Temporal dynamics
		\end{itemize}

		\item \textbf{Speech Transcripts:} Plaintext representations enabling linguistic feature extraction and semantic representation learning
	\end{enumerate}

	\subsection{Data Preprocessing and Normalization}

	\subsubsection{Data Cleaning}

	Preprocessing ensures data quality:

	\begin{itemize}
		\item \textbf{Audio Validation:} Load all WAV files and verify integrity (non-zero length, readable headers)
		\item \textbf{Transcript Cleaning:} Remove XML tags, normalize whitespace, standardize punctuation
		\item \textbf{Missing Data Detection:} Identify subjects with incomplete audio, transcripts, COVAREP, or PHQ-8 labels for exclusion
		\item \textbf{Outlier Identification:} Flag unusual feature values for investigation
	\end{itemize}

	\subsubsection{Feature Normalization Strategy}

	Normalization is applied with careful attention to data leakage prevention:

	\begin{itemize}
		\item \textbf{Per-Subject Normalization:} Compute StandardScaler (zero mean, unit variance) independently for each subject's multimodal features
		\begin{equation}
		\text{x\_normalized} = \frac{\text{x} - \mu_{\text{subject}}}{\sigma_{\text{subject}}}
		\end{equation}

		\item \textbf{Per-Modality Independence:} Audio, text, and acoustic features are normalized separately without cross-modality influence

		\item \textbf{Training/Test Data Separation:} Normalization statistics computed on training data only; applied to validation and test without modification

		\item \textbf{LOSO Consideration:} For each LOSO fold, normalization statistics recomputed using only training subjects to prevent information leakage from test subject
	\end{itemize}

	\subsection{Feature Engineering}

	\subsubsection{Audio Feature Engineering}

	Audio processing maintains temporal structure throughout the extraction pipeline—a design principle critical to capturing depression-related speech dynamics that manifest across time (e.g., reduced prosodic variation, slower speech rate). The system extracts complementary audio features at multiple levels of abstraction:

	\begin{itemize}
		\item \textbf{Self-Supervised Learning (SSL) Embeddings:} Modern SSL models achieve state-of-the-art speech representation through training on massive unlabeled speech corpora. Wav2Vec2 (768-dimensional) is pre-trained on 960 hours of LibriSpeech English speech, learning to predict masked audio segments via a contrastive objective. This produces contextual embeddings capturing phonetic-acoustic patterns without explicit phonetic annotation. WavLM (768-dimensional) extends Wav2Vec 2.0 with multilingual pretraining on 50,000+ hours of speech, providing robust cross-linguistic representations and enhanced stability. Critically, these models output \textit{full temporal sequences}—one embedding vector per 20ms audio frame—preserving all temporal granularity for downstream sequence modeling. No averaging or pooling is performed at extraction time, maintaining the principle from Alhanai et al. that temporal structure must be preserved for depression detection.

		\item \textbf{Traditional Acoustic Features (MFCC):} Mel-Frequency Cepstral Coefficients (13--40 coefficients per frame) capture spectral energy distribution in frequency bands matching human auditory perception. First-order delta features (velocity, $\Delta x_t = x_t - x_{t-1}$) capture frame-to-frame spectral changes reflecting speech rate and articulation. Second-order delta-delta features (acceleration, $\Delta^2 x_t = \Delta x_t - \Delta x_{t-1}$) capture dynamic vocal quality variations. This multi-order expansion yields 39--120 total dimensions per frame while maintaining interpretability: spectral energy changes correlate with emotional prosody, voicing stability, and articulatory control—all potentially affected by depression. The full temporal sequence of MFCC vectors is preserved (T×39-120 tensor).

		\item \textbf{Feature Alignment Strategy:} Audio is segmented into overlapping 3-second chunks (stride 2 seconds) yielding multiple segments per subject. For each chunk, SSL embeddings and MFCC features are extracted at different temporal resolutions (SSL at 20ms, MFCC at 10ms frame rate). Temporal alignment is achieved through bilinear interpolation of MFCC to match SSL resolution, then features are concatenated yielding (T×900) combined audio tensors per chunk. Chunks are aggregated into subject-level sequences by stacking, producing variable-length sequences $T \in [100, 2000+]$ depending on interview duration, with all temporal detail preserved for the hybrid transformer-BiLSTM to exploit.

		\item \textbf{Chunking Strategy Rationale:} Full audio is never processed as a monolithic utterance for three reasons: (1) computational efficiency—3-second chunks fit within GPU memory while maintaining temporal context, (2) numerical stability—SSL models are designed for typical utterance lengths and may exhibit instability on very long sequences, (3) chunk-level aggregation provides implicit data augmentation—different chunk positions within an interview contribute slightly different contexts to the model's learning signal.
	\end{itemize}

	\subsubsection{Text Feature Engineering}

	Linguistic representations extracted from interview transcripts capture semantic and syntactic patterns associated with depression phenomenology. Research in psycholinguistics and computational linguistics has established that depressed individuals exhibit characteristic language patterns: increased first-person singular pronouns (I, me), elevated negative sentiment language, reduced active voice agency markers, and personalization biases in attributional style. The RoBERTa transformer language model, pre-trained on diverse text corpora, learns to recognize these linguistic markers implicitly through distributed representations.

	\begin{itemize}
		\item \textbf{RoBERTa Encoding:} Robustly Optimized BERT Pre-training Approach represents a refinement of the foundational BERT architecture, trained with improved pretraining methodology on larger corpus and longer training schedule. RoBERTa-base (12 layers, 12 attention heads, 768-dimensional embedding space) is applied to interview transcripts through the Hugging Face transformers library. The model processes tokenized transcripts bidirectionally, building contextual understanding where each token's representation depends on surrounding context. For example, the word "I" in ``I feel empty'' receives different representation than identical word in ``I am the team leader,'' enabling discrimination of emotional versus factual contexts. The 768-dimensional output from the model's final layer provides rich semantic representation of entire interview content.

		\item \textbf{CLS Token Pooling Strategy:} Rather than averaging word-level embeddings (which would discard word order and syntactic structure), the special [CLS] token appearing at the beginning of tokenized input is designed during pretraining to aggregate sentence-level semantic meaning. The 768-dimensional [CLS] embedding from the final transformer layer is extracted as the single subject-level text representation. This pooling strategy is theoretically motivated: BERT-family models are explicitly trained such that [CLS] embeddings preserve semantic information necessary for classification tasks, making it a principled universal pooling approach superior to ad-hoc averaging.

		\item \textbf{Transcript Preprocessing:} Minimal text preprocessing is intentionally applied to preserve linguistic markers of depression. Operations include: (1) XML tag removal (speaker labels, timestamps), (2) whitespace normalization (multiple tabs/newlines collapsed to single space), (3) punctuation standardization (no removal, as exclamation marks and dash frequency may be clinically relevant). Aggressive preprocessing commonly employed in NLP (lemmatization, stopword removal) is explicitly \textit{avoided} because such operations discard semantic information: removing first-person pronouns would eliminate depression-relevant markers, and lemmatization conflates "feel" (emotional) with "felt" (temporal aspect), losing clinical nuance.

		\item \textbf{Transcript Source Flexibility:} The system accommodates both manually transcribed high-fidelity transcripts and automatically-generated ASR transcripts via Whisper (OpenAI's robust speech recognition model). While ASR introduces transcription errors, research suggests that semantic information survives moderate ASR error rates (90%+ word accuracy), and RoBERTa's robustness to noise due to pretraining tends to mitigate error impact. This flexibility increases practical deployability by eliminating requirement for expensive manual transcription per interview.
	\end{itemize}

	\subsubsection{Acoustic Feature Engineering}

	COVAREP (Computational Vocal Resonance Period) features represent voice production characteristics reflecting depression-related vocal changes such as reduced pitch variation, decreased speech rate, and altered voice quality. Rather than extracting these features de novo (computationally expensive), the system leverages pre-computed COVAREP matrices available in the DAIC-WOZ dataset.

	\begin{itemize}
		\item \textbf{COVAREP Feature Overview:} The COVAREP algorithm generates frame-level 74-dimensional feature vectors characterizing: (1) fundamental frequency (F0) trajectories including voiced/unvoiced decisions and F0 smoothing coefficients reflective of pitch control, (2) spectral envelope representations from Mel-frequency cepstral coefficients capturing speech energy distribution across frequency bands, (3) voice quality measures including Harmonics-to-Noise Ratio (HNR) reflecting vocal stability and presence of pathological noise, (4) temporal dynamics through delta (velocity) coefficients capturing frame-to-frame changes in vocal characteristics. These 74 dimensions provide rich acoustic characterization developed specifically for speech analysis in clinical contexts.

		\item \textbf{Frame-Level to Subject-Level Aggregation:} When DAIC-WOZ provides frame-level COVAREP features (one 74-dim vector per 10-20ms audio frame), aggregation is necessary to produce fixed-dimensional subject-level representations compatible with the regression architecture. Mean aggregation across time produces a single 74-dim vector capturing average vocal characteristics across the entire interview. This aggregation strategy prioritizes stability and computational efficiency, though alternatives (standard deviation adding temporal variability, percentiles capturing distribution shape) are supported by the architecture and could be explored during hyperparameter optimization.

		\item \textbf{Descriptive Statistics (Optional):} To capture temporal dynamics lost in mean aggregation, the system optionally expands subject-level COVAREP representations to 148 dimensions by concatenating mean and standard deviation of each feature across time. Standard deviations of fundamental frequency, spectral energy, and HNR capture inter-speaker and within-speaker variation in vocal control—higher variance often correlates with more naturalistic emotional expression, while reduced variance may indicate depression-related motor retardation or affective blunting. This expansion trades modest dimensionality increase against preservation of temporal structure within the aggregated representation.
	\end{itemize}

	The multimodal feature approach combining self-supervised learning (SSL) embeddings, RoBERTa language model encodings, and COVAREP acoustic features represents a deliberate architectural choice grounded in comparative analysis against alternative strategies. The comparison evaluates each candidate feature approach along multiple dimensions critical for depression detection: (1) information richness and representation quality relative to depression-related speech characteristics, (2) computational efficiency and practical feasibility given GPU resource constraints, (3) compatibility with temporal modeling requirements—specifically whether temporal structure is preserved throughout processing or prematurely lost to pooling operations (a critical principle established by Alhanai et al. 2018), (4) ability to integrate complementary information from multiple modalities, and (5) alignment with contemporary best practices in speech processing. The proposed SSL+RoBERTa+COVAREP+attention fusion strategy provides optimal balance of performance potential ($R^2 > 0.50$ target), computational tractability, temporal preservation, and multimodal integration. This justifies why approaches such as hand-crafted acoustic features alone (limited capacity), early averaging before sequence modeling (destructive temporal structure loss contradicting Alhanai et al. findings), or simple concatenation fusion (inadequate inter-modal interaction learning) were rejected in favor of the proposed approach.

	\subsection{Model Architecture and Design}

	\subsubsection{Hybrid Transformer-BiLSTM Architecture}

	The proposed model combines transformer and recurrent components with attention-based multimodal fusion:

	\begin{description}
		\item[Input Stage] Audio sequences (B×T×900), optional length information, padding masks

		\item[Transformer Encoder] Captures long-range dependencies and contextual relationships:
		\begin{itemize}
			\item 2 transformer layers
			\item 4--8 multi-head attention heads
			\item 256--384 hidden dimensions
			\item 0.1 dropout for regularization
			\item Key-padding masks handle variable sequence lengths
		\end{itemize}

		\item[BiLSTM Processing] Captures sequential temporal dynamics:
		\begin{itemize}
			\item 1--2 LSTM layers (bidirectional)
			\item 256 hidden state dimension per direction
			\item Output: 512-dim (concatenation of forward/backward)
			\item Processes all timesteps, maintaining temporal resolution
		\end{itemize}

		\item[Attention Pooling] Aggregates temporal information with learned weighting:
		\begin{itemize}
			\item Linear layer projects 512-dim to single attention score per timestep
			\item Softmax normalization (with mask application for padding)
			\item Weighted sum across timesteps yields single 512-dim vector
		\end{itemize}

		\item[Multimodal Fusion] Integrates audio, text, and acoustic modalities:
		\begin{itemize}
			\item Audio: 512-dim pooled representation
			\item Text: 768-dim RoBERTa CLS token
			\item COVAREP: 74-dim acoustic feature vector
			\item Projection to common 256-dim space per modality with layer normalization and ReLU activation
			\item Stack projected vectors into (B×3×256) tensor
			\item Multi-head attention (4 heads) over 3 modality ``tokens''
			\item Attention pooling over modality tokens yields 256-dim fused representation
		\end{itemize}

		\item[Regression Head] Produces continuous PHQ-8 predictions:
		\begin{itemize}
			\item Dense layer: 256→256 with ReLU activation
			\item Dropout: 0.3
			\item Dense layer: 256→64 with ReLU activation
			\item Dense layer: 64→1 producing continuous output
			\item Output range: clamped to [0, 24] representing valid PHQ-8 scores
		\end{itemize}
	\end{description}

	The selection of the hybrid transformer-BiLSTM architecture with attention-based multimodal fusion is justified through systematic comparison with alternative approaches. This analysis evaluates each candidate architecture across critical dimensions: representational capacity for capturing depression-related speech patterns, computational efficiency given the limited dataset size (189 subjects) and available GPU resources, compatibility with temporal sequence modeling requirements, and alignment with published research establishing depression detection through sequence analysis. The hybrid transformer-BiLSTM approach was selected over simpler alternatives (logistic regression, LSTM-only) that lack capacity for sophisticated pattern recognition, and over more complex approaches (pure transformers, large CNNs) that risk severe overfitting on the modest dataset size. This hybrid design balances capacity and regularization, aligning with Alhanai et al. findings.

	\subsubsection{Mathematical Formulation}

	Key mathematical components:

	\begin{align}
	\text{Attention}(Q, K, V) &= \text{softmax}\left(\frac{QK^T}{\sqrt{d_k}} + M\right)V \\
	\text{BiLSTM}_t &= [\text{LSTM}_{\text{fwd}}(t); \text{LSTM}_{\text{back}}(t)] \\
	\alpha_t &= \text{softmax}_{\text{masked}}(w^T \tanh(Wh_t + b)) \\
	v_{\text{pooled}} &= \sum_{t=1}^{T} \alpha_t h_t \\
	\hat{y} &= \text{Linear}_{64 \to 1}(\text{ReLU}(\text{Linear}_{256 \to 64}(\text{ReLU}(\text{Fusion}))))
	\end{align}

	where $M$ denotes attention mask, $d_k$ is head dimension, and Fusion represents multimodal integrated representation.

	\subsection{Training Strategy}

	\subsubsection{LOSO Cross-Validation}

	Leave-One-Subject-Out validation ensures robust generalization:

	\begin{itemize}
		\item For each of $N$ subjects:
		\begin{itemize}
			\item Test set: subject $i$ only
			\item Remaining subjects: partition into validation (10--15\%) and training (85--90\%)
			\item Initialize fresh model (no weight carryover between folds)
			\item Train on training partition
			\item Monitor validation performance
			\item Evaluate once on held-out test subject
		\end{itemize}

		\item Prevents subject-level data leakage ensuring predictions on truly unseen individuals
		\item Provides subject-wise generalization estimates
	\end{itemize}

	\subsubsection{Training Configuration}

	Training proceeds with careful hyperparameter selection:

	\begin{itemize}
		\item \textbf{Loss Function:} Mean Squared Error (MSE), treating all prediction errors equally
		\item \textbf{Optimizer:} AdamW with learning rate $1 \times 10^{-4}$ and weight decay $1 \times 10^{-5}$
		\item \textbf{Batch Size:} 8--16 samples per batch
		\item \textbf{Epoch Budget:} 30--50 epochs per fold
		\item \textbf{Early Stopping:} Patience=5, monitoring validation loss
		\item \textbf{Learning Rate Schedule:} ReduceLROnPlateau with patience=2--3, factor=0.5
		\item \textbf{Gradient Clipping:} Global norm clipping at 1.0 to stabilize training
	\end{itemize}

	\subsubsection{Regularization Strategy}

	Multiple regularization techniques prevent overfitting on limited data:

	\begin{itemize}
		\item \textbf{Dropout:} 0.3 in regression head, 0.1 in transformer layers
		\item \textbf{Early Stopping:} Terminate training when validation loss stalls
		\item \textbf{Weight Decay:} $10^{-5}$ L2 regularization on model parameters
		\item \textbf{Data Augmentation:} Light augmentation during training (noise, pitch variation, time-stretch 0.9--1.1x)
		\item \textbf{Batch Normalization:} Applied implicitly through LayerNorm in transformers
	\end{itemize}

	% =====================================================================
	% 5. FUNCTIONAL REQUIREMENTS
	% =====================================================================

	\section{Functional Requirements}

	\subsection{FR1: Data Input and Validation}

	The system shall accept and validate multimodal interview data:

	\begin{enumerate}
		\item \textbf{FR1.1} Accept WAV format audio files sampled at 16 kHz, mono or stereo configuration
		\item \textbf{FR1.2} Accept plaintext or CSV-formatted speech transcripts with optional speaker labels
		\item \textbf{FR1.3} Accept CSV-formatted COVAREP acoustic feature matrices (74-dimensional)
		\item \textbf{FR1.4} Accept numerical PHQ-8 target labels (0--24 range)
		\item \textbf{FR1.5} Validate file integrity and report errors for corrupted/missing files
		\item \textbf{FR1.6} Generate descriptive error messages guiding data remediation
	\end{enumerate}

	\subsection{FR2: Feature Extraction Pipeline}

	The system shall extract multimodal features preserving temporal structure:

	\begin{enumerate}
		\item \textbf{FR2.1} Extract SSL embeddings (Wav2Vec2, WavLM) maintaining temporal sequences
		\item \textbf{FR2.2} Extract MFCC features with delta and delta-delta coefficients
		\item \textbf{FR2.3} Align and concatenate audio modalities to common temporal resolution
		\item \textbf{FR2.4} Extract RoBERTa embeddings from transcripts
		\item \textbf{FR2.5} Load and aggregate COVAREP acoustic features
		\item \textbf{FR2.6} Perform per-subject feature normalization using StandardScaler
		\item \textbf{FR2.7} Cache all extracted features to disk for efficient reloading
		\item \textbf{FR2.8} Prevent information leakage by computing normalization statistics on training data only
	\end{enumerate}

	\subsection{FR3: Dataset Construction and LOSO Splitting}

	The system shall organize data for rigorous cross-validation:

	\begin{enumerate}
		\item \textbf{FR3.1} Implement LOSO cross-validation with complete subject isolation
		\item \textbf{FR3.2} For each fold, partition data into training (85--90\%), validation (10--15\%), and test (1 subject)
		\item \textbf{FR3.3} Generate PyTorch DataLoader with custom collate functions handling variable-length sequences
		\item \textbf{FR3.4} Apply padding and generate attention masks for variable-length audio
		\item \textbf{FR3.5} Implement stratified splitting considering depression severity distribution
		\item \textbf{FR3.6} Prevent any subject data leakage across train/validation/test splits
	\end{enumerate}

	\subsection{FR4: Model Training and Optimization}

	The system shall train the hybrid neural network architecture:

	\begin{enumerate}
		\item \textbf{FR4.1} Implement transformer encoder (2 layers, 4--8 attention heads, 256--384 hidden dims)
		\item \textbf{FR4.2} Implement bidirectional LSTM (1--2 layers, 256 hidden dim)
		\item \textbf{FR4.3} Implement attention-based temporal pooling over BiLSTM outputs
		\item \textbf{FR4.4} Implement attention-based multimodal fusion combining audio, text, and COVAREP
		\item \textbf{FR4.5} Implement regression head with dense layers and dropout
		\item \textbf{FR4.6} Apply MSE loss function and AdamW optimizer
		\item \textbf{FR4.7} Implement learning rate scheduling with ReduceLROnPlateau
		\item \textbf{FR4.8} Implement early stopping based on validation loss monitoring
		\item \textbf{FR4.9} Apply gradient clipping at norm=1.0
		\item \textbf{FR4.10} Save model checkpoints whenever validation performance improves
	\end{enumerate}

	\subsection{FR5: Model Evaluation and Metrics}

	The system shall compute comprehensive evaluation metrics:

	\begin{enumerate}
		\item \textbf{FR5.1} Compute per-fold $R^2$ (coefficient of determination)
		\begin{equation}
		R^2 = 1 - \frac{\sum (y_i - \hat{y}_i)^2}{\sum (y_i - \bar{y})^2}
		\end{equation}

		\item \textbf{FR5.2} Compute per-fold Mean Absolute Error (MAE)
		\begin{equation}
		\text{MAE} = \frac{1}{n} \sum |y_i - \hat{y}_i|
		\end{equation}

		\item \textbf{FR5.3} Compute per-fold Root Mean Squared Error (RMSE)
		\begin{equation}
		\text{RMSE} = \sqrt{\frac{1}{n} \sum (y_i - \hat{y}_i)^2}
		\end{equation}

		\item \textbf{FR5.4} Aggregate metrics across all LOSO folds (mean, standard deviation)
		\item \textbf{FR5.5} Generate per-subject error analysis
		\item \textbf{FR5.6} Create visualizations: predicted vs. actual scatter plots, residual histograms
		\item \textbf{FR5.7} Output results in JSON and human-readable formats
	\end{enumerate}

	\subsection{FR6: Production Inference}

	The system shall support deployment for real-world inference:

	\begin{enumerate}
		\item \textbf{FR6.1} Provide REST API endpoints accepting audio files and transcripts
		\item \textbf{FR6.2} Perform ASR transcription using Whisper when transcripts not provided
		\item \textbf{FR6.3} Execute complete feature extraction pipeline during inference
		\item \textbf{FR6.4} Load trained model and generate predictions
		\item \textbf{FR6.5} Return structured JSON responses with PHQ-8 prediction and metadata
		\item \textbf{FR6.6} Implement error handling and validation for API requests
		\item \textbf{FR6.7} Support batch inference on multiple subjects
	\end{enumerate}

	% =====================================================================
	% 6. NON-FUNCTIONAL REQUIREMENTS
	% =====================================================================

	\section{Non-Functional Requirements}

	\subsection{NFR1: Performance Requirements}

	The system shall meet the following performance targets:

	\begin{itemize}
		\item \textbf{Prediction Accuracy:} Achieve $R^2 > 0.50$ on test sets using LOSO evaluation
		\item \textbf{Prediction Error:} MAE $< 3.0$ PHQ-8 points, RMSE $< 5.0$ points
		\item \textbf{Training Time:} Complete single LOSO fold training within 2--4 hours on GPU
		\item \textbf{Inference Latency:} Single subject inference (feature extraction + prediction) within 30 seconds on GPU, 2--5 minutes on CPU
		\item \textbf{Throughput:} Batch inference on 100 subjects within 5 minutes on GPU
		\item \textbf{Memory Usage:} Training requires $< 12$ GB GPU VRAM; inference $< 8$ GB
	\end{itemize}

	\subsection{NFR2: Scalability}

	The system architecture shall support scaling to larger datasets:

	\begin{itemize}
		\item \textbf{Dataset Growth:} Support expansion from 189 to 500+ subjects without architectural changes
		\item \textbf{Sequence Length:} Handle interview durations up to 60+ minutes through efficient padding
		\item \textbf{Parallel Processing:} Enable GPU acceleration and distributed inference through API load balancing
		\item \textbf{Feature Caching:} Disk-based caching enables feature reuse across training iterations
	\end{itemize}

	\subsection{NFR3: Reliability and Robustness}

	The system shall operate reliably under expected conditions:

	\begin{itemize}
		\item \textbf{Data Validation:} All inputs validated with descriptive error messages on failure
		\item \textbf{Error Handling:} Graceful exception handling for file I/O, NaN values, GPU memory exhaustion
		\item \textbf{Checkpointing:} Regular model checkpoint saving enables recovery from interruptions
		\item \textbf{Reproducibility:} Fixed random seeds ensure deterministic results across runs
		\item \textbf{Model Validation:} Output predictions clipped to valid PHQ-8 range [0, 24]
		\item \textbf{Data Integrity:} Checksums or file hashing verify cached features integrity
	\end{itemize}

	\subsection{NFR4: Usability}

	The system shall be user-friendly and well-documented:

	\begin{itemize}
		\item \textbf{Configuration:} JSON/YAML configuration files enable hyperparameter adjustment without code modification
		\item \textbf{Command-Line Interface:} CLI commands for training, evaluation, and inference
		\item \textbf{Documentation:} Comprehensive README, API documentation, and example notebooks
		\item \textbf{Logging:} Timestamped logs at multiple verbosity levels for debugging
		\item \textbf{Visualization:} Result plots saved in accessible formats (PNG, PDF)
		\item \textbf{Error Messages:} Clear, actionable error messages guiding users to resolution
	\end{itemize}

	\subsection{NFR5: Security and Privacy}

	The system shall protect sensitive data appropriately:

	\begin{itemize}
		\item \textbf{Data Protection:} Audio files and transcripts not retained after feature extraction
		\item \textbf{Model Security:} Trained model weights stored with restricted file permissions
		\item \textbf{API Authentication:} Optional API key authentication for production deployments
		\item \textbf{Input Sanitization:} API endpoints validate inputs preventing injection attacks and DoS
		\item \textbf{Dependency Management:} Regular security updates for Python libraries
		\item \textbf{Compliance:} Adherence to HIPAA and data protection regulations where applicable
	\end{itemize}

	% =====================================================================
	% 7. EVALUATION & VALIDATION
	% =====================================================================

	\section{Evaluation and Validation}

	\subsection{Evaluation Strategy}

	The system evaluation follows a rigorous protocol:

	\subsubsection{Baseline Comparisons}

	\begin{itemize}
    \item \textbf{BiLSTM:} Captures bidirectional temporal dependencies in sequential audio data, enabling effective modeling of local contextual patterns.

    \item \textbf{Transformer:} Models long-range dependencies using self-attention, allowing the network to focus on relevant temporal segments globally.

    \item \textbf{Single Modality:} Utilizes individual inputs (audio, text, or COVAREP) independently to assess modality-specific contributions.

    \item \textbf{Concatenation Fusion:} Combines features from multiple modalities through direct concatenation without adaptive weighting.

    \item \textbf{Mean Pooling:} Aggregates temporal features by averaging, providing a simple alternative to attention-based pooling.
	\end{itemize}

	\textit{Note: The proposed model combines BiLSTM and Transformer architectures with attention-based multimodal fusion for enhanced representation learning.}

	\subsubsection{Literature Comparison}

	Performance benchmarks against published work:

	\begin{itemize}
		\item Alhanai et al. (2018): $R^2 \approx 0.55$ on PHQ-9 (similar but more difficult task)
		\item Project Target: $R^2 > 0.50$ on PHQ-8 regression is competitive with literature
	\end{itemize}

	\subsection{Critical Analysis}

	\subsubsection{Strengths of the Approach}

	\begin{itemize}
		\item \textbf{Multimodal Integration:} Combines complementary information from audio acoustics, linguistic content, and voice quality
		\item \textbf{Temporal Preservation:} Explicit temporal modeling captures depression-related speech dynamics (prosody, speech rate variation)
		\item \textbf{Transfer Learning:} Pre-trained models (Wav2Vec2, RoBERTa) provide strong feature representations
		\item \textbf{Subject Isolation:} LOSO validation prevents overfitting and provides honest generalization estimates
		\item \textbf{Attention Mechanisms:} Learned inter-modal weighting more principled than fixed-weight combination
		\item \textbf{Regularization:} Dropout, early stopping, and weight decay mitigate overfitting on 189 subjects
	\end{itemize}

	\subsubsection{Limitations and Weaknesses}

	\begin{itemize}
		\item \textbf{Limited Dataset Size:} 189 subjects is modest; larger datasets would enable deeper models
		\item \textbf{Moderate Performance:} $R^2 \approx 0.5$ is acceptable for research but insufficient for clinical deployment
		\item \textbf{Class Imbalance:} Skewed distribution toward mild depression may reduce severe case prediction accuracy
		\item \textbf{Single Dataset Evaluation:} Generalization to other interview formats or populations uncertain
		\item \textbf{Computational Cost:} GPU requirement limits deployment accessibility
		\item \textbf{Temporal Scope:} Single-timepoint assessment; longitudinal tracking not addressed
	\end{itemize}

	\subsection{Expected Outcomes and Performance Scenarios}

	The system's anticipated performance under LOSO cross-validation is presented as a range spanning conservative and optimistic scenarios. The conservative scenario represents realistic expectations given the dataset size, complexity of depression phenotyping, and inherent variability in speech characteristics across individuals. The optimistic scenario reflects performance achievable if contemporary SSL models capture depression-related patterns particularly effectively, and if inter-subject variability in the DAIC-WOZ dataset does not exceed expected bounds. Both scenarios exceed or meet the literature baseline: Alhanai et al. (2018) achieve $R^2 \approx 0.55$ on the more difficult PHQ-9 prediction task (27-point scale) using the same DAIC-WOZ dataset; therefore, our target $R^2 > 0.50$ on PHQ-8 (24-point scale) is meaningfully competitive. Conservative scenario performance ($R^2 = 0.45--0.50$) would demonstrate competitive research-grade results, while optimistic scenario performance ($R^2 = 0.55--0.65$) would significantly exceed published benchmarks, indicating that modern SSL and attention-based approaches indeed provide improvements over prior sequence modeling techniques.

	% =====================================================================
	% 8. IMPLEMENTATION BLUEPRINTS
	% =====================================================================

	\section{Implementation Blueprints}

	This section provides comprehensive architectural blueprints for all three major subsystems of the DepressoSpeech platform: the Frontend user interface, the Backend application server, and the ML Model inference service. Each blueprint includes the technology stack, architectural rationale with comparison tables justifying why each technology was chosen over alternatives, detailed module descriptions, API surface documentation, and the actual project directory structure.

	% ── 8.1 PROJECT STRUCTURE ──────────────────────────────────────────

	\subsection{Project Directory Structure}

	The DepressoSpeech system is organized as a monorepo with three independent services and shared configuration:

	\begin{verbatim}
	DepressoSpeech/
	  |-- Depression-UI/           Frontend (React + Vite)
	  |   |-- src/
	  |   |   |-- components/      Reusable UI components
	  |   |   |-- pages/           Route-level page components
	  |   |   |-- services/        API client layer
	  |   |   |-- hooks/           Custom React hooks
	  |   |   |-- layouts/         Page layout wrappers
	  |   |   |-- utils/           Utility functions
	  |   |   |-- assets/          Static assets
	  |   |   |-- App.jsx          Root component with routing
	  |   |   |-- main.jsx         Application entry point
	  |   |   |-- index.css        Global design system
	  |   |-- public/              Static public assets
	  |   |-- dist/                Production build output
	  |   |-- package.json         NPM dependencies
	  |   |-- vite.config.js       Vite bundler configuration
	  |   |-- tailwind.config.js   Tailwind CSS configuration
	  |
	  |-- backend/                 Backend API (FastAPI)
	  |   |-- src/
	  |   |   |-- routes/          API route handlers
	  |   |   |-- controllers/     Business logic controllers
	  |   |   |-- models/          SQLAlchemy ORM models
	  |   |   |-- services/        Service layer (ML client, email)
	  |   |   |-- middleware/      Auth, metrics middleware
	  |   |   |-- utils/           Utility functions (JWT, hashing)
	  |   |   |-- workers/         Background task workers
	  |   |-- config/              Application settings
	  |   |-- database/            Database initialization
	  |   |-- storage/             Audio file storage
	  |   |-- tests/               Unit and integration tests
	  |   |-- main.py              FastAPI application factory
	  |   |-- requirements.txt     Python dependencies
	  |
	  |-- Model/                   ML Model Service
	  |   |-- src/
	  |   |   |-- api/             FastAPI REST endpoints
	  |   |   |-- inference/       Inference pipeline
	  |   |   |-- model/           Neural network architecture
	  |   |   |-- features/        Feature extraction modules
	  |   |   |-- utils/           Experiment tracker, logging
	  |   |-- scripts/             CLI scripts (train, serve, predict)
	  |   |-- configs/             YAML configuration files
	  |   |-- checkpoints/         Trained model weights
	  |   |-- data/                Dataset files
	  |   |-- logs/                Training and inference logs
	  |   |-- notebooks/           Jupyter analysis notebooks
	  |
	  |-- swagger/                 API Documentation
	  |   |-- index.html           Swagger UI page
	  |   |-- openapi.json         OpenAPI 3.0 specification
	  |   |-- serve.py             Documentation server
	  |
	  |-- SRS.tex                  This document
	  |-- start-all.sh             Linux/macOS startup script
	  |-- start-all.bat            Windows startup script
	\end{verbatim}

	% ── 8.2 SERVICE ARCHITECTURE OVERVIEW ─────────────────────────────

	\subsection{Service Architecture Overview}

	The system follows a microservices architecture with four independently deployable services communicating over HTTP:

	\begin{table}[H]
		\centering
		\caption{Service Port Allocation and Responsibilities}
		\begin{tabular}{|l|c|p{8cm}|}
			\hline
			\textbf{Service} & \textbf{Port} & \textbf{Responsibility} \\ \hline
			Frontend (Vite Dev Server) & 5173 & React single-page application serving the patient, doctor, and admin interfaces \\ \hline
			Backend API (MindScope) & 8000 & Core business logic: authentication, assessments, user management, file storage, dashboard analytics \\ \hline
			ML Model API (DepressoSpeech) & 8001 & Depression severity prediction from audio using trained neural network models \\ \hline
			Swagger Documentation & 8080 & Interactive OpenAPI 3.0 API documentation with Swagger UI \\ \hline
		\end{tabular}
	\end{table}

	\textbf{Data Flow:} The frontend communicates exclusively with the backend via REST API calls (proxied through Vite in development to avoid CORS). The backend orchestrates business logic and delegates ML inference to the Model service. The Swagger server provides unified documentation covering all API endpoints across both backend and ML services.

	% ── 8.3 FRONTEND BLUEPRINT ────────────────────────────────────────

	\subsection{Frontend Blueprint (Depression-UI)}

	\subsubsection{Frontend Technology Stack}

	\begin{itemize}
		\item \textbf{UI Library:} React 19.2 --- component-based, declarative UI with virtual DOM for efficient rendering
		\item \textbf{Build Tool:} Vite 7.3 --- next-generation frontend tooling with native ES module support and instant HMR
		\item \textbf{Routing:} React Router DOM 7.13 --- client-side routing with nested routes and lazy loading
		\item \textbf{Styling:} Tailwind CSS 4.2 --- utility-first CSS framework with JIT compilation
		\item \textbf{Animation:} Framer Motion 11.18 --- production-ready motion library for React
		\item \textbf{Charts:} Recharts 3.8 --- composable charting library built on React and D3
		\item \textbf{Language:} JavaScript (JSX) with ES2020 target
	\end{itemize}

	\subsubsection{Why These Technologies Were Chosen: Frontend Comparison}

	\begin{table}[H]
		\centering
		\caption{Frontend Framework Comparison}
		\small
		\begin{tabular}{|p{2.2cm}|p{2.2cm}|p{2.2cm}|p{2.2cm}|p{2.2cm}|}
			\hline
			\textbf{Criterion} & \textbf{React (Chosen)} & \textbf{Angular} & \textbf{Vue.js} & \textbf{Svelte} \\ \hline
			Learning Curve & Moderate & Steep & Gentle & Gentle \\ \hline
			Ecosystem Size & Largest & Large & Medium & Small \\ \hline
			Component Reuse & Excellent & Good & Good & Limited \\ \hline
			Performance & Very Good (Virtual DOM) & Good & Very Good & Excellent (No VDOM) \\ \hline
			Team Familiarity & High & Low & Medium & Low \\ \hline
			Lazy Loading & Built-in Suspense & Built-in modules & Async components & Limited \\ \hline
			Community Support & Largest & Large & Growing & Small \\ \hline
			Medical UI Libraries & Abundant & Available & Limited & Very Few \\ \hline
		\end{tabular}
	\end{table}

	\textbf{Justification:} React was chosen for its dominant ecosystem, team familiarity, and excellent support for building complex multi-role dashboards (patient, doctor, admin). React's Suspense and lazy loading capabilities enable code splitting across 12 page components, reducing initial bundle size. The component-based architecture maps naturally to the UI requirements: reusable components (Button, Card, Input, Loader) are shared across roles, while role-specific pages (AdminDashboard, DoctorDashboard, Assessment) load on demand.

	\begin{table}[H]
		\centering
		\caption{Build Tool Comparison}
		\small
		\begin{tabular}{|p{2.5cm}|p{2.5cm}|p{2.5cm}|p{2.5cm}|}
			\hline
			\textbf{Criterion} & \textbf{Vite (Chosen)} & \textbf{Webpack (CRA)} & \textbf{Parcel} \\ \hline
			Dev Server Start & < 300ms & 5--30 seconds & 1--5 seconds \\ \hline
			HMR Speed & Instant (<50ms) & 1--3 seconds & 500ms--1s \\ \hline
			Build Speed & Fast (esbuild) & Slow (Babel) & Medium \\ \hline
			Configuration & Minimal & Complex & Zero-config \\ \hline
			ES Module Support & Native & Bundled & Bundled \\ \hline
			Tree Shaking & Excellent (Rollup) & Good & Good \\ \hline
			Proxy Support & Built-in & Built-in & Plugin \\ \hline
		\end{tabular}
	\end{table}

	\textbf{Justification:} Vite provides an order-of-magnitude improvement in developer experience over traditional Webpack-based tools. Its native ES module dev server eliminates bundling during development, delivering sub-second hot module replacement. The built-in proxy configuration seamlessly routes \texttt{/api} requests to the backend (port 8000), avoiding CORS complexity during development.

	\begin{table}[H]
		\centering
		\caption{CSS Framework Comparison}
		\small
		\begin{tabular}{|p{2.5cm}|p{2.5cm}|p{2.5cm}|p{2.5cm}|}
			\hline
			\textbf{Criterion} & \textbf{Tailwind CSS (Chosen)} & \textbf{Bootstrap} & \textbf{Vanilla CSS / CSS Modules} \\ \hline
			Bundle Size & JIT: only used classes & 200+ KB & Manual control \\ \hline
			Customization & Full design system & Limited themes & Unlimited \\ \hline
			Consistency & Enforced via utilities & Component-level & Manual discipline \\ \hline
			Responsive Design & Built-in breakpoints & Grid system & Manual media queries \\ \hline
			Dark Mode & Built-in support & Manual & Manual \\ \hline
			Development Speed & Very Fast & Fast & Slow \\ \hline
		\end{tabular}
	\end{table}

	\textbf{Justification:} Tailwind CSS 4 with JIT compilation produces minimal CSS bundles containing only the utility classes actually used. This is critical for the medical application where consistent, professional styling across patient assessment flows, doctor dashboards, and admin panels must be maintained without custom CSS proliferation. The utility-first approach enables rapid prototyping while maintaining visual consistency.

	\subsubsection{Frontend Component Architecture}

	The frontend implements a role-based single-page application with the following component hierarchy:

	\begin{description}
		\item[Page Components (12 pages):] Each page is lazily loaded via React Suspense to minimize initial bundle size:
		\begin{itemize}
			\item \textbf{Landing} --- Public landing page with feature highlights and call-to-action
			\item \textbf{SignIn / SignUp / VerifyOTP / ForgotPassword} --- Authentication flow with email OTP verification
			\item \textbf{AdminLogin / AdminDashboard} --- System administration with user management, request metrics timeline, ML health monitoring, and severity distribution charts
			\item \textbf{DoctorDashboard} --- Clinical dashboard with patient summaries, high-risk alerts, and score trend visualization
			\item \textbf{Assessment} --- PHQ-8 questionnaire with optional voice recording per question
			\item \textbf{Processing} --- Real-time ML inference progress indicator with stage updates
			\item \textbf{Results} --- Assessment results with PHQ-8 speedometer visualization, confidence intervals, and behavioral feature analysis
			\item \textbf{AssessmentHistory} --- Paginated history of past assessments with trend charts
		\end{itemize}

		\item[Reusable Components (10 components):] Shared across pages:
		\begin{itemize}
			\item \textbf{Navbar} --- Responsive navigation with role-aware menu items and authentication state
			\item \textbf{VoiceRecorder} --- Web Audio API-based recorder with waveform visualization
			\item \textbf{DepessionSpeedometer} --- Animated SVG gauge displaying PHQ-8 severity score
			\item \textbf{ChartPanel / MonitoringTab} --- Recharts-based dashboard visualization components
			\item \textbf{Button / Card / Input / Loader / ResultCard} --- Design system primitives
		\end{itemize}

		\item[Service Layer:] Centralized API client handling JWT token management, request/response interceptors, and error normalization
	\end{description}

	\subsubsection{Frontend Routing and Access Control}

	Route protection is implemented at the component level using role checks:

	\begin{table}[H]
		\centering
		\caption{Frontend Route Access Control Matrix}
		\small
		\begin{tabular}{|l|l|c|c|c|c|}
			\hline
			\textbf{Route} & \textbf{Page} & \textbf{Public} & \textbf{Patient} & \textbf{Doctor} & \textbf{Admin} \\ \hline
			\texttt{/} & Landing & \checkmark & \checkmark & \checkmark & \checkmark \\ \hline
			\texttt{/login} & SignIn & \checkmark & --- & --- & --- \\ \hline
			\texttt{/signup} & SignUp & \checkmark & --- & --- & --- \\ \hline
			\texttt{/verify-otp} & VerifyOTP & \checkmark & --- & --- & --- \\ \hline
			\texttt{/forgot-password} & ForgotPassword & \checkmark & --- & --- & --- \\ \hline
			\texttt{/assessment} & Assessment & --- & \checkmark & --- & --- \\ \hline
			\texttt{/processing} & Processing & --- & \checkmark & --- & --- \\ \hline
			\texttt{/results} & Results & --- & \checkmark & --- & --- \\ \hline
			\texttt{/assessment-history} & AssessmentHistory & --- & \checkmark & --- & --- \\ \hline
			\texttt{/doctor/dashboard} & DoctorDashboard & --- & --- & \checkmark & --- \\ \hline
			\texttt{/admin} & AdminLogin & \checkmark & --- & --- & --- \\ \hline
			\texttt{/admin/dashboard} & AdminDashboard & --- & --- & --- & \checkmark \\ \hline
		\end{tabular}
	\end{table}

	% ── 8.4 BACKEND BLUEPRINT ─────────────────────────────────────────

	\subsection{Backend Blueprint (MindScope API)}

	\subsubsection{Backend Technology Stack}

	\begin{itemize}
		\item \textbf{Framework:} FastAPI 0.100+ --- high-performance async Python web framework with automatic OpenAPI documentation
		\item \textbf{Server:} Uvicorn (ASGI) --- lightning-fast async server based on uvloop and httptools
		\item \textbf{Database:} SQLite with aiosqlite --- embedded database with async driver for zero-configuration deployment
		\item \textbf{ORM:} SQLAlchemy 2.0+ (async mode) --- industry-standard Python ORM with async session support
		\item \textbf{Authentication:} JWT (HS256) via python-jose, bcrypt password hashing
		\item \textbf{Validation:} Pydantic v2 --- data validation using Python type hints
		\item \textbf{Email:} SMTP via smtplib --- OTP delivery for email verification and password reset
		\item \textbf{OAuth:} Google Sign-In (ID token verification)
		\item \textbf{Compression:} GZip middleware for responses over 1 KB
	\end{itemize}

	\subsubsection{Why These Technologies Were Chosen: Backend Comparison}

	\begin{table}[H]
		\centering
		\caption{Backend Framework Comparison}
		\small
		\begin{tabular}{|p{2.2cm}|p{2.2cm}|p{2.2cm}|p{2.2cm}|p{2.2cm}|}
			\hline
			\textbf{Criterion} & \textbf{FastAPI (Chosen)} & \textbf{Flask} & \textbf{Django} & \textbf{Express.js} \\ \hline
			Async Support & Native (ASGI) & Limited (via extensions) & Partial (Django 4+) & Native (Event Loop) \\ \hline
			Auto API Docs & Built-in Swagger & Manual & DRF Browsable & Manual (swagger-jsdoc) \\ \hline
			Type Validation & Pydantic (auto) & Manual / Marshmallow & DRF Serializers & Joi / Zod \\ \hline
			Performance & Very High & Moderate & Moderate & High \\ \hline
			ML Integration & Excellent (Python) & Good (Python) & Good (Python) & Poor (Node.js) \\ \hline
			Learning Curve & Low & Very Low & Moderate & Low \\ \hline
			Background Tasks & Built-in & Celery required & Celery required & Built-in \\ \hline
			WebSocket & Native & flask-socketio & Channels & Native \\ \hline
		\end{tabular}
	\end{table}

	\textbf{Justification:} FastAPI was chosen for three decisive advantages: (1) native async support is essential for handling concurrent audio uploads and ML inference callbacks without blocking; (2) automatic OpenAPI/Swagger documentation generation eliminates documentation drift and enables the unified Swagger UI on port 8080; (3) Python-native execution allows direct integration with the ML model service, sharing data types, serialization formats, and scientific computing libraries. Flask lacks native async support, Django adds unnecessary ORM complexity for a single-database application, and Express.js would introduce a language boundary with the Python ML pipeline, requiring serialization overhead and complicating deployment.

	\begin{table}[H]
		\centering
		\caption{Database Comparison}
		\small
		\begin{tabular}{|p{2.5cm}|p{2.5cm}|p{2.5cm}|p{2.5cm}|}
			\hline
			\textbf{Criterion} & \textbf{SQLite (Chosen)} & \textbf{PostgreSQL} & \textbf{MongoDB} \\ \hline
			Setup Complexity & Zero (file-based) & Server required & Server required \\ \hline
			Deployment & Single file & Docker/install & Docker/install \\ \hline
			Concurrent Reads & Good (WAL mode) & Excellent & Excellent \\ \hline
			Concurrent Writes & Limited & Excellent & Good \\ \hline
			JSON Support & Basic & JSONB (excellent) & Native \\ \hline
			Schema Flexibility & Rigid (relational) & Rigid (relational) & Flexible (document) \\ \hline
			Portability & File copy & Dump/restore & Dump/restore \\ \hline
			Scale Limit & $\sim$100 concurrent users & 10,000+ & 10,000+ \\ \hline
		\end{tabular}
	\end{table}

	\textbf{Justification:} SQLite was chosen because the system is a research tool with moderate concurrency requirements ($\sim$20--50 concurrent users). Its zero-configuration, file-based nature simplifies deployment (no database server to manage), enables trivial backup (file copy), and reduces infrastructure requirements for academic environments. WAL (Write-Ahead Logging) mode provides adequate read concurrency. If the application scales to clinical deployment, migration to PostgreSQL is straightforward via SQLAlchemy's database-agnostic ORM layer.

	\begin{table}[H]
		\centering
		\caption{Authentication Strategy Comparison}
		\small
		\begin{tabular}{|p{2.5cm}|p{2.5cm}|p{2.5cm}|p{2.5cm}|}
			\hline
			\textbf{Criterion} & \textbf{JWT (Chosen)} & \textbf{Session Cookies} & \textbf{OAuth2 Provider} \\ \hline
			Stateless & Yes & No (server-side) & Depends \\ \hline
			Mobile-Friendly & Excellent & Poor (cookie issues) & Good \\ \hline
			Scalability & Horizontal & Requires sticky sessions & Delegated \\ \hline
			Implementation & Moderate & Simple & Complex \\ \hline
			Token Refresh & Refresh tokens & Session renewal & Provider-managed \\ \hline
			Cross-Origin & Header-based & Cookie SameSite & Redirect-based \\ \hline
		\end{tabular}
	\end{table}

	\textbf{Justification:} JWT authentication was chosen for stateless scalability and SPA compatibility. The React frontend stores tokens in memory/localStorage and includes them in Authorization headers, avoiding cross-origin cookie complications. The system implements both access tokens (8-hour TTL) and refresh tokens (7-day TTL) for security. Google OAuth integration provides convenient social login for patients without requiring manual registration.

	\subsubsection{Backend API Architecture}

	The backend API is organized into five route groups under the \texttt{/api/v1} prefix:

	\begin{description}
		\item[Auth Routes (\texttt{/api/v1/auth}):] User registration with email OTP verification, login (email/password), Google OAuth, admin login, password reset, and session management. Implements bcrypt password hashing with salt, 6-digit OTP generation, and configurable expiry (10 minutes default).

		\item[Assessment Routes (\texttt{/api/v1/assessments}):] PHQ-8 questionnaire delivery, assessment creation with per-question answers and optional audio attachments, paginated history listing, latest assessment retrieval, ML processing status polling, and extended ML detail retrieval (confidence intervals, audio quality, behavioral features).

		\item[File Routes (\texttt{/api/v1/files}):] Multipart audio file upload with extension validation (.wav, .mp3, .flac, .ogg, .m4a, .webm), size limit enforcement (100 MB), and local storage management. Returns \texttt{fileId} for association with assessment answers.

		\item[Doctor Routes (\texttt{/api/v1/doctor/dashboard}):] Aggregate analytics (total patients, assessments, high/low risk counts), severity breakdown, recent high-risk alerts with patient details, and per-patient assessment score trends over time.

		\item[Admin Routes (\texttt{/api/v1/admin/dashboard}):] System-wide snapshot (user counts by role, recent registrations, recent assessments), request metrics in 5-minute buckets over 24 hours (latency, error rate), ML model health check (backend and model service status, model load state, inference device).
	\end{description}

	\subsubsection{Backend Data Models}

	The backend persists data using the following SQLAlchemy models:

	\begin{table}[H]
		\centering
		\caption{Backend Database Schema}
		\small
		\begin{tabular}{|l|l|p{7cm}|}
			\hline
			\textbf{Model} & \textbf{Key Fields} & \textbf{Purpose} \\ \hline
			User & id, role, name, email, password\_hash, is\_verified, verification\_otp & Multi-role user accounts with OTP-based email verification \\ \hline
			Assessment & id, user\_id, score\_total, severity, ml\_score, ml\_severity, status & PHQ-8 assessment records with both self-report and ML-predicted scores \\ \hline
			AssessmentAnswer & assessment\_id, question\_id, score, audio\_file\_id & Per-question answer records with optional audio attachment reference \\ \hline
			AssessmentMLDetail & assessment\_id, confidence\_mean, ci\_lower, ci\_upper, audio\_quality\_score & Extended ML prediction metadata \\ \hline
			MediaFile & id, owner\_user\_id, storage\_key, mime\_type, file\_size & Uploaded audio file registry and storage tracking \\ \hline
			RequestMetric & path, method, status\_code, latency\_ms & Per-request performance metrics for admin analytics \\ \hline
		\end{tabular}
	\end{table}

	\subsubsection{Backend-to-ML Communication}

	The backend communicates with the ML model service via an internal REST client (\texttt{MLClient}). When a patient submits an assessment with audio recordings, the backend:

	\begin{enumerate}
		\item Stores the assessment and audio file references synchronously
		\item Triggers a \textbf{background task} (FastAPI's built-in BackgroundTasks) for ML inference
		\item The background task locates the audio file on disk, calls the ML model's \texttt{/predict/extended} endpoint via HTTP
		\item Upon completion, the background task updates the assessment record with \texttt{ml\_score}, \texttt{ml\_severity}, confidence intervals, and audio quality metrics
		\item The frontend polls \texttt{/assessments/\{id\}/processing-status} to track progress
	\end{enumerate}

	This asynchronous design ensures assessment submission returns immediately ($<$100ms) while ML inference (potentially 5--30 seconds) proceeds in the background without blocking the user.

	% ── 8.5 ML MODEL BLUEPRINT ────────────────────────────────────────

	\subsection{ML Model Blueprint (DepressoSpeech API)}

	\subsubsection{ML Model Technology Stack}

	\begin{itemize}
		\item \textbf{Framework:} PyTorch 1.10+ --- dynamic computation graphs enabling flexible model architecture
		\item \textbf{Pre-trained Models:} Hugging Face Transformers 4.15+ --- Wav2Vec2, WavLM, RoBERTa model Hub
		\item \textbf{Audio Processing:} librosa 0.9+, torchaudio 0.10+ --- spectral analysis and waveform manipulation
		\item \textbf{API Server:} FastAPI + Uvicorn --- async HTTP server for inference endpoints
		\item \textbf{Rate Limiting:} SlowAPI --- request throttling for API protection
		\item \textbf{Security:} API key authentication, magic byte validation, file size limits
		\item \textbf{Configuration:} YAML-based inference configuration
		\item \textbf{Experiment Tracking:} Custom SQLite-based experiment tracker
	\end{itemize}

	\subsubsection{Why These Technologies Were Chosen: ML Framework Comparison}

	\begin{table}[H]
		\centering
		\caption{Deep Learning Framework Comparison}
		\small
		\begin{tabular}{|p{2.5cm}|p{2.5cm}|p{2.5cm}|p{2.5cm}|}
			\hline
			\textbf{Criterion} & \textbf{PyTorch (Chosen)} & \textbf{TensorFlow 2.x} & \textbf{JAX} \\ \hline
			Dynamic Graphs & Native (eager mode) & tf.function (mixed) & Functional transforms \\ \hline
			Debugging & Standard Python debugger & Complex (graph mode) & Complex (pure functions) \\ \hline
			Research Adoption & Dominant (80\%+ papers) & Declining in research & Growing but niche \\ \hline
			HuggingFace Integration & First-class & Supported & Limited \\ \hline
			Custom Layers & Pythonic (nn.Module) & Keras Layer API & Flax nn.Module \\ \hline
			Audio Models & Excellent (torchaudio) & tf.audio (limited) & Manual \\ \hline
			Learning Curve & Moderate & Moderate & Steep \\ \hline
			Deployment & TorchServe, ONNX & TF Serving, TFLite & Limited \\ \hline
		\end{tabular}
	\end{table}

	\textbf{Justification:} PyTorch was chosen because the project relies heavily on pre-trained Hugging Face models (Wav2Vec2, WavLM, RoBERTa), which provide first-class PyTorch support. PyTorch's dynamic computation graphs enable rapid prototyping of the hybrid transformer-BiLSTM architecture with attention-based fusion. Its dominance in research (80\%+ of recent NLP/speech papers use PyTorch) ensures community support and reproducibility against published baselines.

	\begin{table}[H]
		\centering
		\caption{Audio Feature Extraction Comparison}
		\small
		\begin{tabular}{|p{2.2cm}|p{2.5cm}|p{2.5cm}|p{2.5cm}|}
			\hline
			\textbf{Criterion} & \textbf{SSL + MFCC + COVAREP (Chosen)} & \textbf{Hand-Crafted Only (eGeMAPS)} & \textbf{End-to-End Raw Audio} \\ \hline
			Feature Quality & State-of-the-art & Good but limited & Depends on data volume \\ \hline
			Data Requirement & Low (transfer learning) & Low & Very High ($>$10K hours) \\ \hline
			Interpretability & Moderate & High & Very Low \\ \hline
			Computation Cost & High (GPU for SSL) & Low (CPU) & Very High \\ \hline
			Temporal Structure & Fully preserved & Depends on pipeline & Preserved \\ \hline
			Depression-Specific & Multi-perspective & Voice-quality focused & Learned (opaque) \\ \hline
			Modality Coverage & Audio + Text + Acoustic & Audio only & Audio only \\ \hline
		\end{tabular}
	\end{table}

	\textbf{Justification:} The multimodal SSL + MFCC + COVAREP approach was chosen because: (1) Self-supervised learning embeddings (Wav2Vec2, WavLM) capture sophisticated acoustic-phonetic patterns learned from 50,000+ hours of unlabeled speech, enabling strong representations despite only 189 labeled subjects; (2) MFCC features provide interpretable spectral characteristics that are well-established depression biomarkers; (3) COVAREP acoustic features capture voice quality parameters (F0, HNR, shimmer, jitter) directly linked to depression symptomology in clinical literature; (4) RoBERTa text embeddings capture semantic depression markers (negative affect, first-person pronouns) that are independent of acoustic quality. End-to-end raw audio processing was rejected because the limited dataset size (189 subjects) is insufficient for learning meaningful representations from scratch.

	\begin{table}[H]
		\centering
		\caption{Model Architecture Comparison}
		\small
		\begin{tabular}{|p{2.2cm}|p{2.5cm}|p{2.5cm}|p{2.5cm}|}
			\hline
			\textbf{Criterion} & \textbf{Transformer + BiLSTM (Chosen)} & \textbf{Pure Transformer} & \textbf{BiLSTM Only} \\ \hline
			Parameters & $\sim$5.2M & $\sim$8--15M & $\sim$2--3M \\ \hline
			Long-Range Deps & Excellent (attention) & Excellent & Moderate (vanishing gradient) \\ \hline
			Local Patterns & Excellent (LSTM) & Moderate & Excellent \\ \hline
			Overfitting Risk & Balanced & High (too many params) & Low (limited capacity) \\ \hline
			Training Stability & Good & Moderate (warmup needed) & Good \\ \hline
			Attention Fusion & Multi-head over modalities & Native & Concatenation only \\ \hline
			Param/Subject Ratio & $\sim$27.5K/subject & $\sim$42--79K/subject & $\sim$10--16K/subject \\ \hline
		\end{tabular}
	\end{table}

	\textbf{Justification:} The hybrid transformer-BiLSTM architecture ($\sim$5.2M parameters) strikes a deliberate balance between capacity and regularization for the 189-subject dataset. The transformer encoder (2 layers, 4--8 heads) captures long-range dependencies across the interview (e.g., consistent monotone patterns spanning minutes), while the BiLSTM captures local temporal dynamics (pause duration, speech rate variation within utterances). This complementary design achieves an approximate 27,500 parameters per training subject---well within the empirical safe zone for moderate datasets. A pure transformer would exceed 8M parameters, risking severe overfitting with a parameter-to-subject ratio exceeding 42,000:1. A BiLSTM-only model lacks the capacity to capture global interview-level patterns crucial for depression detection.

	\begin{table}[H]
		\centering
		\caption{Multimodal Fusion Strategy Comparison}
		\small
		\begin{tabular}{|p{2.2cm}|p{2.5cm}|p{2.5cm}|p{2.5cm}|}
			\hline
			\textbf{Criterion} & \textbf{Attention Fusion (Chosen)} & \textbf{Concatenation} & \textbf{Late Fusion (Averaging)} \\ \hline
			Inter-Modal Learning & Learned weighting & None (static) & None (equal weight) \\ \hline
			Adaptivity & Dynamic per sample & Fixed & Fixed \\ \hline
			Parameters Added & $\sim$0.9M & $\sim$0.1M & 0 \\ \hline
			Missing Modality & Graceful (reweight) & Zeroed dimensions & Dropped average \\ \hline
			Interpretability & Attention weights viewable & N/A & N/A \\ \hline
			Published Efficacy & State-of-the-art & Common baseline & Simple baseline \\ \hline
		\end{tabular}
	\end{table}

	\textbf{Justification:} Attention-based multimodal fusion (4-head attention over 3 modality tokens projected to 256-dim) enables the model to dynamically weight modalities per subject. For clear-speech subjects, the model can up-weight text semantic features; for subjects with distinctive vocal patterns but normal language, audio/acoustic features receive higher attention. This adaptive weighting is superior to concatenation (which treats all modalities equally regardless of signal quality) and late fusion (which cannot learn cross-modal interactions). The attention weights also provide interpretability: clinicians can inspect which modalities contributed most to each prediction.

	\subsubsection{ML API Endpoints}

	The ML model service exposes the following REST endpoints:

	\begin{table}[H]
		\centering
		\caption{ML Model API Endpoints}
		\small
		\begin{tabular}{|l|l|l|p{5cm}|}
			\hline
			\textbf{Method} & \textbf{Endpoint} & \textbf{Rate Limit} & \textbf{Description} \\ \hline
			POST & \texttt{/predict} & 30/min & Single audio file prediction returning PHQ-8 score and severity \\ \hline
			POST & \texttt{/predict/extended} & 20/min & Extended prediction with confidence intervals, audio quality, and behavioral features \\ \hline
			POST & \texttt{/predict/batch} & 10/min & Batch prediction for up to 20 audio files \\ \hline
			GET & \texttt{/health} & Unlimited & Health check returning model load status and device \\ \hline
		\end{tabular}
	\end{table}

	\subsubsection{ML Security Measures}

	The ML API implements defense-in-depth security:

	\begin{enumerate}
		\item \textbf{API Key Authentication (SEC-2):} Optional \texttt{X-API-Key} header validation using constant-time comparison (\texttt{secrets.compare\_digest})
		\item \textbf{Rate Limiting (SEC-3):} Per-endpoint rate limits via SlowAPI (30/min predict, 10/min batch)
		\item \textbf{File Size Limits:} Maximum 100 MB per uploaded file
		\item \textbf{Content Validation (SEC-8):} Magic byte verification ensures uploaded files are actual audio (checks RIFF, ID3, fLaC, OggS, WebM headers)
		\item \textbf{Input Sanitization:} Participant ID validated against regex pattern (alphanumeric, 1--64 chars)
		\item \textbf{Thread Safety (BP-1):} Pipeline lock prevents concurrent model access, ensuring prediction consistency
		\item \textbf{TLS Advisory (SEC-5):} Warning logged when SSL is disabled, reminding operators about HIPAA compliance
	\end{enumerate}

	% ── 8.6 CROSS-VALIDATION IMPLEMENTATION ───────────────────────────

	\subsection{Cross-Validation and Training Implementation}

	\begin{enumerate}
		\item \textbf{Data Preparation:} Load DAIC-WOZ files, validate integrity, exclude subjects with missing modalities

		\item \textbf{Feature Extraction:} Run parallel feature extraction pipeline for all subjects, cache results to disk as PyTorch tensors

		\item \textbf{LOSO Training Loop:}
		\begin{enumerate}
			\item For each subject $S$ (fold):
			\begin{enumerate}
				\item Create train/val/test splits with $S$ in test
				\item Initialize new model with random weights
				\item Compute per-modality normalization on training split only
				\item Train for 30--50 epochs with early stopping (patience=5)
				\item Apply ReduceLROnPlateau scheduling (patience=2--3, factor=0.5)
				\item Checkpoint best model by validation $R^2$
				\item Evaluate on held-out subject $S$
				\item Store per-fold metrics ($R^2$, MAE, RMSE)
			\end{enumerate}
		\end{enumerate}

		\item \textbf{Results Aggregation:} Compute cross-fold mean and standard deviation of metrics

		\item \textbf{API Deployment:} Package best model checkpoint, deploy inference pipeline via FastAPI on port 8001

		\item \textbf{Documentation:} Generate Swagger UI on port 8080 via unified OpenAPI specification

		\item \textbf{Visualization:} Generate scatter plots, residual histograms, and per-subject error analysis
	\end{enumerate}

	% ── 8.7 DEPLOYMENT AND STARTUP ────────────────────────────────────

	\subsection{Deployment and Startup}

	\subsubsection{Unified Startup}

	All four services are started simultaneously using the provided startup scripts:

	\begin{itemize}
		\item \textbf{Linux/macOS:} \texttt{./start-all.sh} --- Bash script with colored output, dependency checking, port conflict resolution, and process management with Ctrl+C graceful shutdown
		\item \textbf{Windows:} \texttt{start-all.bat} --- Batch script launching services in separate minimized windows with \texttt{taskkill}-based cleanup
	\end{itemize}

	Both scripts support flags: \texttt{--no-deps} (skip dependency installation) and \texttt{--kill-only} (release occupied ports without starting services).

	\subsubsection{API Documentation (Swagger)}

	Unified API documentation is served on port 8080 using an OpenAPI 3.0 specification covering both the Backend and ML Model service endpoints. The Swagger UI has been custom-designed with a professional dark-theme interface featuring:

	\begin{itemize}
		\item \textbf{Custom Dark Theme:} Purpose-built CSS overriding the default Swagger UI with a professional dark color scheme (\texttt{\#0f172a} background, \texttt{\#6366f1} primary accents), optimized typography using Inter and JetBrains Mono fonts, and generous spacing for readability
		\item \textbf{Branded Header:} Sticky header with project branding, version badge (v2.0.0), OpenAPI 3.0 badge, and service count indicator with animated status dots
		\item \textbf{Service Port Strip:} Visual reference bar showing all four service ports (Frontend :5173, Backend :8000, ML Model :8001, Docs :8080) for quick reference during development
		\item \textbf{Color-Coded Methods:} HTTP methods distinguished by left-border color (green for GET, blue for POST, orange for PUT, red for DELETE) with subtle background tinting for visual scanning
		\item \textbf{Interactive ``Try it out'' functionality} for testing endpoints directly from the browser
		\item \textbf{JWT Bearer authentication and API Key authentication support} via the Authorize dialog
		\item \textbf{Complete request/response schema documentation} with examples and syntax-highlighted JSON
		\item \textbf{Endpoint grouping by service tag:} Auth, Assessments, Files, Doctor Dashboard, Admin Dashboard, ML Prediction, ML Health --- totaling 24+ documented endpoints
		\item \textbf{Responsive Design:} Layout adapts gracefully to mobile and tablet viewports with adjusted padding and stacked navigation
	\end{itemize}

	% ── 8.8 PERFORMANCE ENGINEERING ───────────────────────────────────

	\subsection{Performance Engineering}

	System-wide performance optimization is critical for a responsive user experience across the assessment flow (audio recording $\rightarrow$ upload $\rightarrow$ ML inference $\rightarrow$ results display). This subsection documents all latency-reduction measures implemented across the frontend, backend, and infrastructure layers.

	\subsubsection{Frontend Performance Optimizations}

	The React frontend implements multiple strategies to minimize perceived and actual latency:

	\begin{table}[H]
		\centering
		\caption{Frontend Performance Optimization Inventory}
		\small
		\begin{tabular}{|p{3.5cm}|p{4cm}|p{4.5cm}|}
			\hline
			\textbf{Optimization} & \textbf{Mechanism} & \textbf{Impact} \\ \hline
			CORS Preflight Elimination & API base URL changed from absolute \texttt{http://localhost:8000} to proxy-relative \texttt{/api/v1}; Vite dev server proxies to backend & Removes 100--300ms CORS OPTIONS preflight from \textit{every} cross-origin request \\ \hline
			In-Memory Response Cache & GET responses cached in \texttt{Map} with 15-second TTL; stale entries auto-evicted & Eliminates redundant network calls when components re-mount rapidly (e.g., tab switches) \\ \hline
			In-Flight Request Dedup & Concurrent identical GET requests share a single \texttt{Promise} & Prevents N parallel fetches when N components request the same data simultaneously \\ \hline
			Parallel Dashboard Fetch & Doctor dashboard's 3 sequential API calls (\texttt{summary}, \texttt{trends}, \texttt{alerts}) replaced with \texttt{Promise.all} & Dashboard load time reduced from $3 \times RTT$ to $1 \times RTT$ ($\sim$3$\times$ faster) \\ \hline
			Aggressive Code Splitting & Vendor (React, Router), Charts (Recharts), and Animation (Framer Motion) in separate chunks & Initial JS bundle reduced; returning visits leverage browser HTTP cache per chunk \\ \hline
			Vite Module Warmup & \texttt{server.warmup.clientFiles} pre-transforms \texttt{App.jsx}, \texttt{Landing.jsx}, \texttt{Navbar.jsx}, \texttt{api.js} at startup & First page load avoids on-demand transformation latency \\ \hline
			Dependency Pre-Bundling & Heavy libraries (\texttt{framer-motion}, \texttt{recharts}) added to \texttt{optimizeDeps.include} & Browser receives pre-bundled ESM; no runtime parsing overhead \\ \hline
			Lazy Loading with Suspense & All 12 page components loaded via \texttt{React.lazy()} with \texttt{Suspense} fallback & Only active route's code downloaded; idle routes deferred \\ \hline
		\end{tabular}
	\end{table}

	\subsubsection{Backend Performance Optimizations}

	The FastAPI backend implements both per-request and architectural optimizations:

	\begin{table}[H]
		\centering
		\caption{Backend Performance Optimization Inventory}
		\small
		\begin{tabular}{|p{3.5cm}|p{4cm}|p{4.5cm}|}
			\hline
			\textbf{Optimization} & \textbf{Mechanism} & \textbf{Impact} \\ \hline
			Batched Metrics Writes & \texttt{MetricsMiddleware} queues \texttt{RequestMetric} objects in a \texttt{deque} buffer; background \texttt{asyncio.Task} flushes batch every 5 seconds via single \texttt{session.add\_all()} & Eliminates per-request DB write from hot path, saving 5--20ms per API call \\ \hline
			Aggressive GZip Compression & \texttt{GZipMiddleware} threshold lowered from 1024 to 500 bytes & More responses compressed; typical JSON payloads see 60--80\% size reduction \\ \hline
			SQLite WAL Mode & \texttt{PRAGMA journal\_mode=WAL} enables concurrent readers with single writer & Reads no longer block on writes; critical for polling endpoints like \texttt{processing-status} \\ \hline
			SQLite Page Cache & \texttt{PRAGMA cache\_size=-32000} (32 MB) and \texttt{temp\_store=MEMORY} & Frequently accessed table pages served from RAM; reduces disk I/O \\ \hline
			Connection Pool & SQLAlchemy \texttt{pool\_size=10}, \texttt{max\_overflow=20}, \texttt{pool\_pre\_ping=True} & Connections reused across requests; no per-request TCP/TLS handshake overhead \\ \hline
			Async ML Inference & Assessment submission triggers \texttt{BackgroundTasks} for ML prediction; frontend polls \texttt{/processing-status} & Submission returns in $<$100ms; ML inference (5--30s) runs asynchronously \\ \hline
			Skipped Metrics Paths & Health, favicon, docs, and OpenAPI endpoints excluded from metrics tracking & Eliminates unnecessary buffer entries for high-frequency non-business endpoints \\ \hline
		\end{tabular}
	\end{table}

	\subsubsection{Performance Comparison: Before vs. After Optimization}

	\begin{table}[H]
		\centering
		\caption{Measured Latency Impact of Optimizations}
		\small
		\begin{tabular}{|p{4cm}|c|c|c|}
			\hline
			\textbf{Operation} & \textbf{Before (ms)} & \textbf{After (ms)} & \textbf{Improvement} \\ \hline
			Typical GET API call (frontend $\rightarrow$ backend) & 150--400 & 20--80 & 3--5$\times$ faster \\ \hline
			Doctor dashboard full load & 900--1200 & 300--400 & $\sim$3$\times$ faster \\ \hline
			Cached GET (identical re-request within 15s) & 150--400 & $<$1 & Instant (memory) \\ \hline
			Backend per-request overhead (metrics write) & 5--20 & $<$0.1 & $\sim$100$\times$ faster \\ \hline
			Initial page load (Vite HMR cold start) & 800--1500 & 300--500 & 2--3$\times$ faster \\ \hline
			JSON payload transfer (typical 4 KB response) & Full size & 60--80\% smaller & GZip at 500B threshold \\ \hline
		\end{tabular}
	\end{table}

	\subsubsection{Caching Strategy Comparison}

	The frontend caching approach was selected through systematic evaluation:

	\begin{table}[H]
		\centering
		\caption{Frontend Caching Strategy Comparison}
		\small
		\begin{tabular}{|p{2.5cm}|p{2.5cm}|p{2.5cm}|p{2.5cm}|}
			\hline
			\textbf{Criterion} & \textbf{In-Memory Cache (Chosen)} & \textbf{Service Worker Cache} & \textbf{React Query / SWR} \\ \hline
			Implementation & 40 lines vanilla JS & Complex lifecycle & Library dependency \\ \hline
			Bundle Impact & Zero bytes & SW registration & 10--30 KB gzipped \\ \hline
			Cache Invalidation & Automatic on mutations & Manual cache busting & Built-in revalidation \\ \hline
			Offline Support & No & Yes & Partial \\ \hline
			Request Dedup & Custom (included) & Not built-in & Built-in \\ \hline
			Complexity & Minimal & High & Moderate \\ \hline
		\end{tabular}
	\end{table}

	\textbf{Justification:} The lightweight in-memory cache was chosen because: (1) zero bundle size impact aligns with the code-splitting strategy; (2) 15-second TTL is appropriate for a medical assessment application where data freshness is important but sub-second polling is unnecessary; (3) automatic invalidation on POST/PUT/DELETE mutations ensures users always see updated data after actions; (4) in-flight request deduplication is included natively, whereas Service Workers would require additional logic. React Query or SWR would add unnecessary library weight for an application with straightforward data-fetching patterns.

	% =====================================================================
	% 9. FUTURE ENHANCEMENTS
	% =====================================================================

	\section{Future Enhancements and Research Directions}

	\subsection{Proposed Improvements}

	\begin{itemize}
		\item \textbf{Multi-Task Learning:} Jointly predict PHQ-8 score, anxiety (GAD-7), and sleep quality
		\item \textbf{Uncertainty Quantification:} Bayesian neural networks or ensemble methods generating prediction intervals
		\item \textbf{Domain Adaptation:} Transfer learning to other interview datasets (e.g., different interviewer, setting)
		\item \textbf{Explainability:} Attention visualization and LIME-based feature attribution
		\item \textbf{Longitudinal Modeling:} Time-series models tracking depression trajectory across multiple interviews
		\item \textbf{Real-Time Streaming:} Process audio segments in real-time during interviews
	\end{itemize}

	\subsection{Research Opportunities}

	\begin{itemize}
		\item Cross-linguistic evaluation on non-English interviews
		\item Investigation of speech biomarkers predictive of treatment response
		\item Federated learning enabling multi-site training without data centralization
		\item Integration with wearable biosensors (heart rate variability, sleep data)
		\item Clinical validation through prospective studies comparing predictions to clinician diagnoses
	\end{itemize}

	% =====================================================================
	% 10. CONCLUSION
	% =====================================================================

	\section{Conclusion}

	This Software Requirements Specification defines a comprehensive full-stack system for speech-based depression severity prediction using multimodal machine learning. The system comprises three independently deployable services: (1) a React-based frontend providing role-specific interfaces for patients (PHQ-8 assessment with voice recording), doctors (clinical dashboards with trend analysis), and administrators (system monitoring with ML health tracking); (2) a FastAPI backend (MindScope) handling authentication, data management, and ML inference orchestration; and (3) a dedicated ML Model API (DepressoSpeech) implementing the hybrid transformer-BiLSTM architecture with attention-based multimodal fusion for PHQ-8 score prediction from audio. All services are documented through a unified OpenAPI 3.0 Swagger interface served on port 8080, featuring a custom-designed dark-theme UI with professional typography, color-coded HTTP methods, and responsive layout.

	The specifications detail functional requirements (feature extraction, model training, evaluation, REST API endpoints) and non-functional requirements (performance, scalability, security). Every technology choice---from React and Vite for the frontend, FastAPI and SQLite for the backend, to PyTorch and Hugging Face for the ML pipeline---is justified through systematic comparison tables evaluating alternatives along critical dimensions including performance, team familiarity, ecosystem maturity, and alignment with project constraints. Implementation is grounded in the comprehensive blueprints presented in Section~8, ensuring the final system meets all stated requirements while positioned for future enhancements and clinical validation.

	Version 2.1 introduces systematic performance engineering across the entire stack (Section~8.8). Frontend optimizations---including CORS preflight elimination via Vite proxy, in-memory response caching with 15-second TTL, in-flight request deduplication, and parallelized dashboard API calls---reduce typical API round-trip latency by 3--5$\times$. Backend optimizations---including batched metrics database writes (eliminating per-request DB overhead), aggressive GZip compression at 500-byte threshold, SQLite WAL mode with 32~MB page cache, and connection pooling---further reduce server-side processing time. Together, these optimizations ensure a responsive user experience where assessment submissions return in under 100ms, dashboard loads complete in under 400ms, and repeated data fetches resolve in under 1ms from cache.

	Success is measured by achieving $R^2 > 0.50$ regression performance with MAE $< 3.0$ PHQ-8 points across all LOSO folds, establishing a competitive baseline for speech-based mental health assessment at the intersection of computational linguistics, acoustic analysis, and machine learning. The complete system is deployable via a single startup script (\texttt{./start-all.sh} or \texttt{start-all.bat}) that initializes all four services with dependency management, port conflict resolution, and graceful shutdown handling.

	\newpage

	% =====================================================================
	% APPENDICES
	% =====================================================================

	\section*{Appendix A: Configuration Parameters}

	Key hyperparameters and their meanings:

	\begin{itemize}
		\item \textbf{transformer\_layers = 2}: Number of stacked transformer encoder layers for capturing long-range dependencies.

		\item \textbf{transformer\_heads = 4}: Multi-head attention heads within each transformer layer for parallel attention mechanisms.

		\item \textbf{transformer\_hidden = 256}: Hidden dimension size in transformer layers; larger values increase capacity but require more GPU memory.

		\item \textbf{lstm\_layers = 2}: Number of bidirectional LSTM layers for sequential temporal pattern modeling.

		\item \textbf{lstm\_hidden = 256}: Hidden state dimension for LSTM cells; concatenated with reverse direction yielding 512-dim output.

		\item \textbf{fusion\_dim = 256}: Common projection dimension for multimodal fusion; all modalities projected to this space before attention.

		\item \textbf{dropout = 0.3}: Regularization probability; 30\% of neurons randomly disabled during training to prevent overfitting.

		\item \textbf{optimizer = AdamW}: Adaptive moment estimation with weight decay; combines momentum with L2 regularization for stable convergence.

		\item \textbf{learning\_rate = 1e-4}: Initial learning rate for optimizer; controls step size during gradient updates.

		\item \textbf{weight\_decay = 1e-5}: L2 regularization coefficient penalizing large weights to reduce overfitting on 189-subject dataset.

		\item \textbf{batch\_size = 16}: Number of samples per gradient update; small batch stabilizes training on limited data.

		\item \textbf{epochs = 50}: Maximum training iterations over full training set before early stopping.

		\item \textbf{early\_stopping\_patience = 5}: Number of validation epochs without improvement before halting to prevent overfitting.

		\item \textbf{grad\_clip\_norm = 1.0}: Maximum L2 norm for gradient clipping; prevents exploding gradients during backpropagation.

		\item \textbf{audio\_chunk\_size = 3}: Duration in seconds for audio processing chunks; balances temporal context with computational efficiency.

		\item \textbf{audio\_stride = 2}: Overlap between consecutive chunks in seconds; smaller stride increases temporal resolution.

		\item \textbf{sample\_rate = 16000}: Audio sampling frequency in Hz; 16 kHz captures speech frequencies up to 8 kHz per Nyquist theorem.

		\item \textbf{mfcc\_n\_mels = 40}: Number of mel-frequency cepstral coefficient bands; standard for speech processing.

		\item \textbf{loso\_val\_split = 0.15}: Fraction of training set reserved for per-fold validation; 15\% for validation, 85\% for training.
	\end{itemize}

	\section*{Appendix B: Mathematical Formulations}

	\subsection*{Multimodal Attention Fusion}

	\begin{align}
	Q, K, V &\in \mathbb{R}^{B \times 3 \times 64} \\
	\text{Score} &= QK^T / \sqrt{64} \\
	\text{Attention} &= \text{softmax}(\text{Score}) \cdot V \\
	\text{Output} &= \text{Concat}(\text{Attention}_1, \ldots, \text{Attention}_4) \cdot W^O
	\end{align}

	\subsection*{Regression Loss}

	\begin{equation}
	\mathcal{L} = \frac{1}{B} \sum_{b=1}^{B} (\hat{y}_b - y_b)^2
	\end{equation}

	\section*{Appendix C: Common Errors and Troubleshooting}

	\begin{itemize}
		\item \textbf{GPU Out of Memory:} Reduce batch size or sequence chunk size
		\item \textbf{NaN Loss:} Check for data normalization; verify target values in [0, 24]
		\item \textbf{Poor Convergence:} Try higher learning rate or reduce gradient clipping threshold
		\item \textbf{Data Leakage Errors:} Verify normalization statistics computed on training data only
		\item \textbf{Port Already in Use:} Run \texttt{./start-all.sh --kill-only} to release all ports before restarting
		\item \textbf{CORS Errors in Browser:} Ensure API calls use relative path \texttt{/api/v1} (not absolute \texttt{http://localhost:8000}) so the Vite proxy handles them
		\item \textbf{Slow Dashboard Load:} Verify the in-memory API cache is active; check browser DevTools Network tab for duplicate requests
		\item \textbf{SQLite Lock Errors:} Confirm WAL mode is enabled (\texttt{PRAGMA journal\_mode} should return \texttt{wal})
		\item \textbf{Stale Data After Mutation:} The frontend cache auto-invalidates on POST/PUT/DELETE; if stale data persists, call \texttt{invalidateCache()} manually
	\end{itemize}

	\section*{Appendix D: Performance Optimization Inventory}

	Complete list of all performance optimizations implemented in version 2.1:

	\subsection*{D.1 Frontend Optimizations (7 items)}

	\begin{enumerate}
		\item \textbf{CORS Preflight Elimination:} API base URL changed from absolute \texttt{http://localhost:8000/api/v1} to relative \texttt{/api/v1}. Vite's built-in proxy (\texttt{server.proxy.\/api $\rightarrow$ http://localhost:8000}) forwards requests server-side, appearing as same-origin to the browser. This eliminates the preflight OPTIONS request that browsers send for cross-origin requests with custom headers (\texttt{Authorization}, \texttt{Content-Type: application/json}), saving 100--300ms per API call.

		\item \textbf{In-Memory Response Cache:} A lightweight \texttt{Map}-based cache (\texttt{\_cache}) stores GET response data with 15-second time-to-live (TTL). Cache keys combine HTTP method, path, and authorization header hash to ensure per-user isolation. On mutation requests (POST, PUT, DELETE), the entire cache is cleared to guarantee data freshness.

		\item \textbf{In-Flight Request Deduplication:} A second \texttt{Map} (\texttt{\_inflight}) tracks pending GET requests by cache key. If an identical request is initiated while another is in flight, the existing \texttt{Promise} is returned instead of creating a new network request. This coalesces burst traffic common in React's Strict Mode (double-render in development) and concurrent component mounts.

		\item \textbf{Parallel Dashboard Fetching:} The doctor dashboard's \texttt{getDashboardSnapshot()} function previously issued three sequential API calls (summary, patient-trends, alerts). These are now wrapped in \texttt{Promise.all()}, executing concurrently. Total latency drops from $3 \times \text{RTT}$ to $1 \times \text{RTT}$.

		\item \textbf{Vite Module Warmup:} The \texttt{server.warmup.clientFiles} configuration pre-transforms frequently accessed modules (\texttt{App.jsx}, \texttt{Landing.jsx}, \texttt{Navbar.jsx}, \texttt{api.js}) during dev server startup. This shifts transformation cost from first-request time to startup time, reducing perceived latency for the initial page load.

		\item \textbf{Granular Chunk Splitting:} \texttt{rollupOptions.output.manualChunks} separates the vendor bundle into three independently cacheable chunks: \texttt{vendor} (React, ReactDOM, React Router), \texttt{charts} (Recharts), and \texttt{motion} (Framer Motion). App code changes don't invalidate heavy library caches.

		\item \textbf{Aggressive Dependency Pre-Bundling:} \texttt{optimizeDeps.include} expanded from 3 to 5 entries, adding \texttt{framer-motion} and \texttt{recharts}. Vite pre-bundles these into optimized ESM at startup, eliminating runtime transformation overhead during navigation.
	\end{enumerate}

	\subsection*{D.2 Backend Optimizations (4 items)}

	\begin{enumerate}
		\item \textbf{Batched Metrics Middleware:} The \texttt{MetricsMiddleware} previously called \texttt{session.add()} and \texttt{session.commit()} on every request, adding 5--20ms of SQLite write latency to the critical path. The optimized version appends \texttt{RequestMetric} objects to a bounded \texttt{deque} (max 500 entries) and flushes the batch every 5 seconds via a background \texttt{asyncio.Task} using a single \texttt{session.add\_all()} + \texttt{commit()}, reducing per-request overhead to $<$0.1ms.

		\item \textbf{GZip Compression Threshold:} \texttt{GZipMiddleware} minimum size reduced from 1024 to 500 bytes. Typical API responses (assessment lists, dashboard snapshots, processing status) range 600--4000 bytes; the lower threshold ensures these are compressed, achieving 60--80\% size reduction and proportionally faster network transfer.

		\item \textbf{SQLite Performance Pragmas:} Four SQLite pragmas applied at connection time: \texttt{journal\_mode=WAL} (concurrent reads), \texttt{synchronous=NORMAL} (balanced durability/speed), \texttt{cache\_size=-32000} (32~MB page cache), \texttt{temp\_store=MEMORY} (temporary tables in RAM). Combined impact: read-heavy workloads (dashboard polling) execute with minimal disk I/O.

		\item \textbf{Connection Pool Tuning:} SQLAlchemy async engine configured with \texttt{pool\_size=10}, \texttt{max\_overflow=20}, and \texttt{pool\_pre\_ping=True}. This maintains 10 persistent connections (expandable to 30 under load) and validates connections before use, eliminating per-request TCP handshake and authentication overhead.
	\end{enumerate}

\end{document}
